% \iffalse meta-comment
%
% Copyright (C) 2026 by Pierre Senellart
%
% This work may be distributed and/or modified under the
% conditions of the LaTeX Project Public License, either version 1.3
% of this license or (at your option) any later version.
% The latest version of this license is in
%   http://www.latex-project.org/lppl.txt
% and version 1.3 or later is part of all distributions of LaTeX
% version 2005/12/01 or later.
%
% This work has the LPPL maintenance status `maintained'.
%
% The Current Maintainer of this work is Pierre Senellart
% <pierre@senellart.com> and a version control system for this work
% is available at http://github.com/PierreSenellart/proofgraph
%
% This work consists of the files proofgraph.dtx and proofgraph.ins
% and the derived file proofgraph.sty.
%
% \fi
%
% \iffalse
%<*driver>
\documentclass{ltxdoc}
\usepackage{hypdoc}
\usepackage{tikz}
\usepackage[autorun,cite,citelabel=tag]{proofgraph}
% This documentation is itself a proofgraph document: the example results in
% the introduction are typeset with the package enabled, and the graph shown
% there is the one they produce (rendered automatically thanks to autorun).
% The cite option is on, so the \cite in the example proof becomes a node too.
% TikZ and hyperref (via hypdoc) are loaded, so the nodes of that graph are
% clickable: each links to the result it represents.
\newtheorem{theorem}{Theorem}
\newtheorem{lemma}{Lemma}
\newtheorem{corollary}{Corollary}
\proofgraphstyle{theorem}{shape=box,style=filled,fillcolor=gold}
\proofgraphstyle{corollary}{shape=box,style="rounded,filled",fillcolor=palegoldenrod}
\proofgraphstyle{lemma}{shape=ellipse,style=filled,fillcolor=lightgray}
\makeatletter
% Verb-free renderers for the command-syntax labels of the description lists in
% the Usage section (\verb / |...| is not allowed inside \item[...]).
\newcommand\cmdname[1]{{\normalfont\ttfamily\char`\\#1}}
\newcommand\cmdarg[1]{{\normalfont\ttfamily\char`\{}\meta{#1}%
  {\normalfont\ttfamily\char`\}}}
\newcommand\optname[1]{{\normalfont\ttfamily#1}}
\newcommand\cmdoarg[1]{{\normalfont\ttfamily[}\meta{#1}{\normalfont\ttfamily]}}
% The code examples are set as verbatim, sometimes wrapped in quote for
% indentation; give every display a small, balanced vertical margin (about half
% a line). The outermost environment governs the spacing: a quote-wrapped
% verbatim takes it from \quote (the inner verbatim is the first list item, whose
% top \topsep is suppressed), a bare verbatim from \verbatim itself, so both are
% set to the same \topsep. \quote is a \list, which re-applies \@listi after any
% \pretocmd, so its spacing goes inside the \list parameter argument; \verbatim
% is a \trivlist, where a \pretocmd suffices. A verbatim display leaves more
% space below than above (the return to the surrounding baselineskip), so a small
% negative skip after every verbatim trims the bottom to match the top.
\newskip\pgphvbskip
\AtBeginDocument{\pgphvbskip=.5\baselineskip}
\renewcommand\quote{\list{}{\rightmargin\leftmargin
  \topsep\pgphvbskip\partopsep\z@\parsep\z@}\item\relax}
\pretocmd{\@verbatim}{\topsep\pgphvbskip\partopsep\z@}{}{}
\AfterEndEnvironment{verbatim}{\vspace{-11.5\p@}}
\makeatother
% Long unbreakable inline code spans (|\proofgraphedge{...}|, URLs) would
% otherwise stick into the margin; let TeX loosen lines enough to rebreak.
\setlength\emergencystretch{3em}
\EnableCrossrefs
\CodelineIndex
% Two index columns rather than three, so long control-sequence names
% (\@spopargbegintheorem, \pgfmathsetlengthmacro) fit without overflowing.
\setcounter{IndexColumns}{2}
\RecordChanges
\begin{document}
  \DocInput{proofgraph.dtx}
\end{document}
%</driver>
% \fi
%
% \CheckSum{2262}
%
% \CharacterTable
%  {Upper-case    \A\B\C\D\E\F\G\H\I\J\K\L\M\N\O\P\Q\R\S\T\U\V\W\X\Y\Z
%   Lower-case    \a\b\c\d\e\f\g\h\i\j\k\l\m\n\o\p\q\r\s\t\u\v\w\x\y\z
%   Digits        \0\1\2\3\4\5\6\7\8\9
%   Exclamation   \!     Double quote  \"     Hash (number) \#
%   Dollar        \$     Percent       \%     Ampersand     \&
%   Acute accent  \'     Left paren    \(     Right paren   \)
%   Asterisk      \*     Plus          \+     Comma         \,
%   Minus         \-     Point         \.     Solidus       \/
%   Colon         \:     Semicolon     \;     Less than     \<
%   Equals        \=     Greater than  \>     Question mark \?
%   Commercial at \@     Left bracket  \[     Backslash     \\
%   Right bracket \]     Circumflex    \^     Underscore    \_
%   Grave accent  \`     Left brace    \{     Vertical bar  \|
%   Right brace   \}     Tilde         \~}
%
% \changes{v0.1.0}{2026/05/31}{Initial version}
% \changes{v1.0.0}{2026/06/02}{First released version}
% \changes{v1.0.1}{2026/07/09}{Fix an infinite loop when, with the
%   \texttt{cite} option, a citation occurred in a proof nested within
%   another proof: the capture no longer reinstalls over itself}
% \changes{v1.0.1}{2026/07/10}{Support, with the \texttt{cite} option,
%   the two-optional-argument citation form
%   \texttt{[pre][post]\{key\}} of natbib and biblatex; it
%   previously broke compilation with confusing errors inside proofs}
% \changes{v1.0.1}{2026/07/10}{Capture \texttt{citet} and \texttt{citep}
%   (including starred variants) like \texttt{cite} when they are defined}
% \changes{v1.1.0}{2026/07/27}{With \texttt{autorun}, run Graphviz only when
%   the graph, the engine or the format has changed. Rendering at every
%   compilation kept producing a different image, since Graphviz stamps a
%   creation date into its \textsf{cairo}-based outputs, and
%   \textsf{latexmk} would then rerun until it reported that too many passes
%   were needed}
% \changes{v1.1.0}{2026/07/27}{Capture, with the \texttt{cite} option, every
%   citation command of \textsf{natbib} and \textsf{biblatex}, such as
%   \texttt{textcite} and \texttt{parencite}, down to those printing a mere
%   fragment of a citation such as \texttt{citeauthor}, and add
%   \texttt{proofgraphcitecommand} to capture further ones}
% \changes{v1.1.0}{2026/07/27}{Require a LaTeX2e format of October 2020 or
%   later, which the use of \textsf{expl3} and \texttt{NewDocumentCommand}
%   already implied}
% \changes{v1.1.0}{2026/07/27}{Capture \texttt{zcref} of
%   \textsf{zref-clever}, in its starred, optional-argument and multiple-label
%   forms, like the other cross-referencing commands}
% \changes{v1.1.0}{2026/07/27}{Emit citation nodes again when
%   \texttt{citelabel!=tag} finds no tag to use. A stray \texttt{iftrue}, counted
%   as a conditional while a false conditional was being skipped over, made the
%   skip overshoot and swallow the writes emitting the node, which Graphviz then
%   drew from its edges alone: unstyled, labelled by its internal name and
%   without a hyperlink. It affected every first run, and \textsf{biblatex}
%   documents always}
% \changes{v1.1.0}{2026/07/27}{Recover the citation tag of the numeric and
%   alphabetic \textsf{biblatex} styles for \texttt{citelabel!=tag}, which
%   records no citation tag in the \texttt{.aux} for it to read}
% \changes{v1.1.0}{2026/07/28}{Add \texttt{proofgraphproofenv}, declaring
%   further environments to be treated as proofs, for a document whose proofs
%   do not live in an environment named \texttt{proof}}
% \changes{v1.1.0}{2026/07/28}{Capture cross-references and citations in the
%   body of a result and not only in its proof and its optional title, a
%   statement resting on an earlier result being a dependency like any other;
%   the new \texttt{statements} option, on by default, turns this off}
% \changes{v1.1.0}{2026/07/28}{Add \texttt{usescite} and
%   \texttt{proofgraphciteedge}, recording a dependency on cited work by hand,
%   for a proof that does not cite it visibly and for a result stated without a
%   proof to hold a citation}
% \changes{v1.1.0}{2026/07/28}{Capture citations, and not only
%   cross-references, in the optional title of a result}
% \changes{v1.1.0}{2026/07/28}{Trim the spaces after the commas of the
%   lists given to \texttt{proofgraphtrack} and \texttt{proofgraphuntrack};
%   every name but the first was silently ignored}
% \changes{v1.1.0}{2026/07/29}{Keep the braces of an optional title that is
%   a single brace group, the bracing the \textsf{amsthm} manual prescribes to
%   hide the \texttt{]} of a citation note. A title such as
%   \texttt{[\{cite[p.~37]\{a-l\}\}]} reached the heading unbraced, whose own
%   scan for \texttt{]} then stopped inside the note: the citation vanished from
%   the printed title and two errors were reported against the capture}
% \changes{v1.1.0}{2026/07/29}{Let \texttt{proofof} pin a proof as a whole,
%   wherever in the proof it stands: the references made before it went to
%   whichever result preceded the proof, which a proof deferred to an appendix
%   made easy to hit. The result it names may now be labelled further on in the
%   document}
% \changes{v1.1.0}{2026/07/29}{Point the hyperlink of a citation node at the
%   destination \textsf{biblatex} gives a bibliography entry, which carries the
%   number of the reference section before the key, rather than at the one the
%   LaTeX kernel and \textsf{natbib} give, which \textsf{biblatex} never
%   defines: the node led to the end of the document instead of the entry}
% \changes{v1.1.0}{2026/07/29}{Set the note of a citation node inside its
%   tag and after it, ``[Knu68, pp.~23--27]'', the way a citation with a note is
%   printed, and render its \texttt{--} and \texttt{---} as the dashes the
%   document shows}
% \changes{v1.1.0}{2026/07/30}{Let a \texttt{proofof} in a proof reach the
%   proofs nested in it, the steps of a structured proof: a nested proof took as
%   its result the one that preceded the proof around it, so the references it
%   made escaped the pin. A nested proof stated after a claim still proves that
%   claim}
% \changes{v1.1.0}{2026/07/30}{Record no dependency while the page is
%   assembled: a cross-reference or a citation in a sectional title is executed
%   again in the running head, and a page breaking inside a proof made it a
%   dependency of the result that proof belonged to}
% \changes{v1.1.0}{2026/07/30}{Draw a result stated inside a proof, a claim
%   settled in the course of an argument, as supporting the result that proof
%   establishes. It was drawn with its own proof below it but nothing above,
%   detached from the result it was raised to establish}
% \changes{v1.1.0}{2026/07/30}{Report an editing rule that dropped nothing,
%   \texttt{proofgraphignore} above all, whose two labels go the way the
%   dependency is recorded, as in \texttt{proofgraphedge}, and not the way the
%   arrow is drawn: a rule written the wrong way round changed nothing and said
%   nothing. The warning names the rule to write instead. Say so in the manual
%   as well, which described the arguments in terms of the edge drawn}
% \changes{v1.1.0}{2026/07/30}{Add \texttt{proofgraphignorecite}, dropping
%   the dependency of a result on a cited work, which \texttt{proofgraphignore}
%   cannot name, both of its arguments being labels}
% \changes{v1.1.0-dev}{2026/08/03}{Run Graphviz through \texttt{ShellEscape}
%   of \textsf{shellesc} rather than through the \texttt{write18} primitive,
%   which under LuaTeX writes to an ordinary output stream instead of running
%   anything: \texttt{autorun} did nothing under LuaLaTeX, whatever
%   shell-escape was passed, unless the document happened to load
%   \textsf{shellesc}, as one loading \textsf{minted} does}
% \changes{v1.1.0-dev}{2026/08/03}{Report a rendering that produced no image,
%   naming what stood in the way: shell escape disabled, or restricted and
%   therefore not permitting Graphviz, or Graphviz missing. The run went on as
%   if it had rendered, and the only sign was the warning of
%   \texttt{proofgraph} on the next run, saying that the image was not there}
% \changes{v1.1.0-dev}{2026/08/03}{Say, when \texttt{proofgraph} finds no
%   image to embed and \texttt{autorun} is on, that the image is rendered at
%   the end of the run and a further run embeds it, rather than asking for
%   shell escape and the option the document already has}
% \changes{v1.1.0-dev}{2026/08/03}{Read the attributes given to
%   \texttt{proofgraphstyle} and \texttt{proofgraphstylecite} with \texttt{\#}
%   as an ordinary character, so that a Graphviz colour may be written
%   \texttt{"\#ffd700"}. It stopped the run with ``Illegal parameter number'',
%   leaving HSV as the only way to a colour of one's own}
% \changes{v1.1.0-dev}{2026/08/04}{Read the plain file of Graphviz with
%   \texttt{\#} as an ordinary character too: its node lines echo the style
%   attributes, so an RGB colour reached the run that embeds the hyperlinked
%   graph and stopped it with three further ``Illegal parameter number''
%   errors, whether or not the colour was still in the source}
% \changes{v1.1.0-dev}{2026/09/02}{Capture, under \textsf{apxproof}, the
%   proofs it typesets through its saved copy of the proof environment. A
%   plain proof, whether left in the body or belonging to a
%   \texttt{toappendix} block and replayed with it, contributed no edges at
%   all; only the deferred proof of a repeated theorem, and a proof whose
%   title names its result, were captured}
% \changes{v1.1.0-dev}{2026/09/02}{Let the label passed by the
%   \textsf{apxproof} hook pin the proof instance the wrapping has opened,
%   rather than replace the current proof. A \texttt{proofof} in a deferred
%   proof redirected the current proof globally, the redirection outlived
%   the proof, and the edges of what followed shifted to the named result}
% \changes{v1.1.0-dev}{2026/09/02}{Scan the title of a proof with
%   \texttt{protect} at its typesetting meaning, not emptied. Under
%   \textsf{babel} (so under \textsf{lipics}, which loads it), a tie in a
%   title such as ``Proof of Theorem\texttt{\textasciitilde}\texttt{ref...}''
%   expanded through \texttt{active@prefix}, whose fallback reproduces the
%   bare active character when \texttt{protect} is empty, and the run died
%   of a full input stack}
%
% \DoNotIndex{\newcommand,\renewcommand,\providecommand,\def,\edef,\gdef,\xdef}
% \DoNotIndex{\let,\global,\long,\relax,\expandafter,\csname,\endcsname}
% \DoNotIndex{\ifcsname,\ifdefined,\ifx,\fi,\else,\@ifstar,\@ifnextchar}
% \DoNotIndex{\begingroup,\endgroup,\@empty,\@ne,\advance,\the,\string}
% \DoNotIndex{\immediate,\write,\newwrite,\openout,\closeout,\space}
% \DoNotIndex{\RequirePackage,\ProvidesPackage,\NeedsTeXFormat}
% \DoNotIndex{\g@addto@macro,\@for,\do,\@gobble,\@firstofone}
% \DoNotIndex{\ifcase,\or,\@makeother}
%
% \let\embedproofgraph\proofgraph
% \def\proofgraph{\textsf{proofgraph}}
% \def\apxproof{\textsf{apxproof}}
% \def\Graphviz{\textsf{Graphviz}}
% \def\LuaTeX{Lua\TeX}
% \def\LuaLaTeX{Lua\LaTeX}
%
% \GetFileInfo{proofgraph.sty}
%
% \title{The \proofgraph\ package}
% \author{Pierre Senellart \\ \texttt{pierre@senellart.com}}
% \date{\filedate \quad \fileversion}
%
% \maketitle
%
% \begin{abstract}
%   The \proofgraph\ package automatically produces a graph of the
%   dependencies between the results (theorems, lemmas, propositions)
%   of a mathematical article. It infers an edge from one result to
%   another whenever the proof of the former refers to the latter
%   through an ordinary cross-reference, so that most dependencies need
%   no manual annotation; commands are provided for those a visible
%   reference does not capture. The graph is written as a \Graphviz\
%   |.dot| file, which can be rendered externally or, optionally,
%   embedded into the document automatically.
% \end{abstract}
%
% \section{Introduction}
%
% When writing or reading a mathematical article, it is often useful to visualise
% how the results depend on one another: which lemma is used in the
% proof of which theorem, and so on. \proofgraph\ builds such a
% dependency graph automatically. Each theorem-like environment becomes
% a node; each cross-reference (|\ref|, |\cref|\ldots) appearing
% \emph{inside a proof} becomes an edge from the proved result to the
% referred-to result. Usually no manual annotation is required, and commands
% are provided for the dependencies a visible reference does not capture.
%
% The package is best explained by a small example of its use. This manual is
% itself a \proofgraph\ document, so the results stated below are typeset
% with the package enabled, and the graph shown at the end of this section is
% the one they produce.
%
% Load \proofgraph\ in the preamble, before the |\newtheorem| commands that
% declare your theorem-like environments (so it sees them as they are created).
% Here, as in this manual, we pass three options: |autorun| renders the graph
% with \Graphviz\ and embeds it automatically (it needs shell-escape); |cite|
% captures the |\cite|s made inside proofs as dependencies on external work; and
% |citelabel=tag| labels those citation nodes by their printed tag:
% \begin{quote}
% \begin{verbatim}
% \usepackage[autorun,cite,citelabel=tag]{proofgraph}
% \newtheorem{theorem}{Theorem}
% \newtheorem{lemma}{Lemma}
% \newtheorem{corollary}{Corollary}
% \end{verbatim}
% \end{quote}
% Then state and prove, as usual, a few results that build on one another:
% \begin{quote}
% \begin{verbatim}
% \begin{lemma}\label{lem:even}
%   The sum of two even integers is even.
% \end{lemma}
% \begin{lemma}\label{lem:odd}
%   The sum of two odd integers is even.
% \end{lemma}
% \begin{theorem}\label{thm:parity}
%   The sum of two integers of the same parity is even.
% \end{theorem}
% \begin{proof}
%   By Lemmas~\ref{lem:even} and~\ref{lem:odd}; see~\cite{euclid}.
% \end{proof}
% \begin{corollary}\label{cor:double}
%   For every integer $n$, the number $2n$ is even.
% \end{corollary}
% \begin{proof}
%   Apply Theorem~\ref{thm:parity} to $n$ and $n$.
% \end{proof}
% \end{verbatim}
% \end{quote}
% We obtain the following results, exactly as without the package:
% \begin{lemma}\label{lem:even}
%   The sum of two even integers is even.
% \end{lemma}
% \begin{lemma}\label{lem:odd}
%   The sum of two odd integers is even.
% \end{lemma}
% \begin{theorem}\label{thm:parity}
%   The sum of two integers of the same parity is even.
% \end{theorem}
% \begin{proof}
%   By Lemmas~\ref{lem:even} and~\ref{lem:odd}; see~\cite{euclid}.
% \end{proof}
% \begin{corollary}\label{cor:double}
%   For every integer $n$, the number $2n$ is even.
% \end{corollary}
% \begin{proof}
%   Apply Theorem~\ref{thm:parity} to $n$ and $n$.
% \end{proof}
% In addition, \proofgraph\ writes out the dependency graph. With the |autorun|
% option it is rendered with \Graphviz\ and included wherever you call
% |\proofgraph|:
% \begin{quote}
% \begin{verbatim}
% \proofgraph[width=0.7\textwidth]
% \end{verbatim}
% \end{quote}
% which here produces:
% \begin{center}
% \embedproofgraph[width=0.7\textwidth]
% \end{center}
% Each cross-reference made inside a proof becomes an arrow \emph{into} the
% result that proof establishes: the two lemmas point to
% Theorem~\ref{thm:parity}, which in turn points to Corollary~\ref{cor:double}.
% The |\cite| in the proof of Theorem~\ref{thm:parity} is captured in the same
% way, thanks to the |cite| option: the external work appears as a distinct
% node (drawn note-shaped, and labelled here by its printed citation tag,
% ``\texttt{[1]}'', because of |citelabel=tag|) with an arrow into the theorem
% it helps prove.
% Moreover, since this manual loads \textsf{hyperref} and \textsf{TikZ}, the
% graph's nodes are \emph{clickable}: try clicking one to jump to the
% result it represents (see |\proofgraph| and the |hyperlinks| option).
%
% A handful of commands tune such a graph: |\proofgraphstyle| sets the
% appearance of each kind of node (the colours here come from it); |\uses| and
% |\proofgraphedge| record a dependency that no visible |\ref| expresses;
% |\proofof| attaches a proof to its result; |\proofgraphuntrack| keeps an
% environment out of the graph; and |\proofgraphignore| and |\proofgraphexclude|
% drop an unwanted edge or node. The next section puts several of these to work
% on a real article; Section~\ref{sec:usage} then documents them, and the
% package options, in full.
%
% \section{A real-world example}
% \label{sec:realworld}
% The introductory example is deliberately tiny. To show what \proofgraph\
% produces on a genuine article, the graph in Figure~\ref{fig:realworld} is the
% one obtained from \emph{Connecting Knowledge Compilation Classes and Width
% Parameters}~\cite{AmarilliCMS20}. It has 62 numbered results, drawn
% from five environments (|result|, |theorem|, |proposition|, |corollary| and
% |lemma|), depends on 24 external works, and contains 114
% edges in all.
%
% \begin{figure}[p]
% \centering
% \makebox[\linewidth]{%
%   \includegraphics[width=\dimexpr\linewidth+2cm\relax]{example-realworld}}
% \caption{The graph \proofgraph\ produces for~\cite{AmarilliCMS20}. The colours
%   and shapes make the
%   structure legible, with the orange double-bordered introduction
%   \texttt{result}s on the right, fed by the gold theorems, the blue and
%   pale-gold propositions and corollaries, and the grey lemmas.}
% \label{fig:realworld}
% \end{figure}
%
% The graph was produced from the paper essentially as published. Beyond loading
% the package, the only additions were the following.
% \begin{itemize}
% \item The |cite| option was switched on, and citation nodes labelled by their
%   printed tag, by loading the package as
%   |\usepackage[cite=true,citelabel=tag]{proofgraph}|. The |\cite|s made inside
%   proofs then appear as the note-shaped nodes (``\texttt{[17]}'',
%   ``\texttt{[46, Theorem 4.10]}''\dots) along the left of the graph; a work
%   cited for two different results, or for two of its own theorems, gives one
%   node per |\cite| note, so those 23~works appear as 39~nodes.
% \item One environment was kept out of the graph with a single
%   |\proofgraphuntrack{observation}| (a lone preliminary observation that
%   nothing else uses).
% \item The five result kinds were styled, with decreasing prominence, by five
%   |\proofgraphstyle| declarations:
% \begin{quote}
% \begin{verbatim}
% \proofgraphstyle{result}{shape=box,style="filled,bold",%
%                          fillcolor=orange,peripheries=2}
% \proofgraphstyle{theorem}{shape=box,style=filled,%
%                           fillcolor=gold,penwidth=2}
% \proofgraphstyle{corollary}{shape=box,style="rounded,filled",%
%                             fillcolor=palegoldenrod}
% \proofgraphstyle{proposition}{shape=box,style="rounded,filled",%
%                               fillcolor=lightblue}
% \proofgraphstyle{lemma}{shape=ellipse,style=filled,%
%                         fillcolor=lightgray}
% \end{verbatim}
% \end{quote}
% \item Eleven |\proofgraphedge| commands, recording 27 edges in total, were
%   added for the dependencies that no |\ref| inside a |proof| expresses:
%   eight ``obvious'' corollaries stated without a |proof| environment
%   (|\proofgraphedge{cor:upper_bound}{thm:upper_bound,...}|), and three theorems
%   whose proof is developed across a whole section rather than inside a single
%   |proof|.
% \item One |\proofgraphciteedge|, for a lemma stated without a proof, its proof
%   being in the conference version of the article.
% \end{itemize}
% Nothing else was annotated: every other edge, and the numbering of all 62
% nodes, was inferred automatically from the ordinary cross-references the proofs
% already contained. In particular, no |\uses|, |\proofof|, |\proofgraphignore|
% or |\proofgraphexclude| was needed. The five |result| environments that restate
% the headline results in the introduction
% (|\begin{result}[(Corollary~\ref{cor:obdd_omega})]|\dots) are linked to the
% body results they preview through the |\ref| in their optional title, captured
% automatically (see the note on reference-carrying titles in
% Section~\ref{sec:usage}). Two compilation runs (plus \BibTeX, needed once for
% the |tag| labels) suffice; the |.dot| file is then rendered with \Graphviz,
% |dot -Tpdf main.dot -o main.pdf|, as with any other document.
%
% \section{Usage}
% \label{sec:usage}
%
% \subsection{Default behaviour}
% Load \proofgraph\ in the preamble. It only needs to come \emph{after} whatever
% provides your theorem environments, that is, after |\usepackage{amsthm}| if you
% load it yourself (some classes, such as \textsf{lipics} or \textsf{acmart},
% provide theorems for you, in which case there is nothing to load first), and
% \emph{before} the |\newtheorem| commands that declare your theorem-like
% environments, so that \proofgraph\ sees them as they are created:
% \begin{verbatim}
%   \usepackage{amsthm}      % or nothing, if your class provides theorems
%   \usepackage{proofgraph}
%   \newtheorem{theorem}{Theorem}
%   \newtheorem{lemma}{Lemma}
% \end{verbatim}
% Then write your document as usual. After two compilation runs, a file
% |\jobname.dot| is produced. Render it with \Graphviz:
% \begin{verbatim}
%   dot -Tpdf myfile.dot -o mygraph.pdf
% \end{verbatim}
% Every \emph{numbered} result of a theorem-like environment you declare becomes
% a node, and every cross-reference (|\ref|, |\cref|\dots) inside a proof becomes
% an edge into the result that proof establishes, with no annotation. The
% commands captured are |\ref|, |\cref| and |\Cref| (\textsf{cleveref}),
% |\autoref| (\textsf{hyperref}), |\eqref| (\textsf{amsmath}), |\vref|
% (\textsf{varioref}) and |\zcref| (\textsf{zref-clever}, in every form:
% starred, with options, and with a comma-separated list of labels), each of
% them whenever it is defined. Purely page-related references (|\pageref|,
% |\cpageref|, |\zcpageref|) are not captured, as they point at a place rather
% than at a result; |\uses| records a dependency they would have expressed.
% Nodes are
% labelled by the result's formatted name and number (``Theorem~1''), not the
% environment name. Unnumbered (starred) results are skipped, as they are usually
% expository or repeated statements; |\proofgraphtrack| draws a chosen one anyway.
% A cross-reference in the \emph{optional title} of a result is captured the same
% way:
% |\begin{theorem}[(generalises~\ref{thm:x})]| in a result labelled \meta{A} adds
% an edge from \meta{A} to |thm:x|, handy for restatements (e.g., an introduction
% that previews results with |\begin{result}[(Corollary~\ref{cor:y})]|). The
% \emph{body} of a result is read the same way: a statement that refers to an
% earlier result, ``Under the hypotheses of Theorem~2\dots'', depends on it just
% as a proof would. Set |statements=false| to build the graph from proofs alone.
%
% A result \emph{stated inside a proof}, a claim raised and settled in the course
% of an argument, is drawn as a node of its own and given an edge to the result
% that proof establishes: the claim is a step of that proof, and the graph says
% so without any annotation. Its own proof, whether it follows the claim as a
% nested |proof| environment or as the rest of the enclosing one, hangs below the
% claim, so the argument appears in the graph with the shape it has on the page.
% An unnumbered claim, like any unnumbered result, is left out, and its
% dependencies with it; number it, or name it in |\proofgraphtrack|, to have it
% drawn.
%
% \subsection{Recording dependencies by hand}
% \begin{description}
% \item[\cmdname{uses}\cmdarg{label}]\SpecialUsageIndex{\uses}%
%   Inside a proof, records a dependency on the result labelled \meta{label} even
%   when no visible reference is made. A comma-separated list is accepted, as in
%   |\uses{lem:a,lem:b}|.
% \item[\cmdname{proofgraphedge}\cmdarg{A}\cmdarg{labels}]%
%   \SpecialUsageIndex{\proofgraphedge}%
%   Adds edges recording that result \meta{A} depends on the result(s)
%   \meta{labels} (a comma-separated list), like a |\uses| but stated outside any
%   proof. Handy for an ``obvious'' corollary that relies on earlier results
%   without a |proof| environment, as in |\proofgraphedge{cor:x}{thm:y,lem:z}|. It
%   may appear anywhere; all labels are resolved when the graph is written.
% \item[\cmdname{usescite}\cmdoarg{note}\cmdarg{keys}]%
%   \SpecialUsageIndex{\usescite}%
%   Inside a proof, the |\uses| of external work: records a dependency on the
%   bibliography entries \meta{keys} (a comma-separated list) and draws them as
%   citation nodes, without citing them visibly. The optional \meta{note} is the
%   one |\cite| would have taken, as in |\usescite[Thm~3]{knuth}|; it labels the
%   node and, as with a captured |\cite|, tells one result of a work from
%   another.
% \item[\cmdname{proofgraphciteedge}\cmdoarg{note}\cmdarg{A}\cmdarg{keys}]%
%   \SpecialUsageIndex{\proofgraphciteedge}%
%   The |\proofgraphedge| of citation nodes: records that result \meta{A}
%   depends on the cited work(s) \meta{keys}, stated outside any proof, as in
%   |\proofgraphciteedge{thm:x}{knuth}|. Useful for a result stated without a
%   proof, which has no place to hold a citation. Like |\proofgraphedge| it may
%   appear anywhere, all labels being resolved when the graph is written.
%
%   Both commands create the citation node when the work has none yet, and both
%   work whether or not the |cite| option is on: that option only says whether a
%   plain |\cite| is captured on its own. A work reached both ways, once by
%   |\cite| capture and once by hand, still yields a single node.
% \item[\cmdname{proofof}\cmdarg{label}]\SpecialUsageIndex{\proofof}%
%   By default a proof is attached to the most recently started result. A
%   detached proof whose optional title names its result is attached
%   automatically: |\begin{proof}| with |[of Theorem~\ref{thm:x}]| attaches the
%   proof to |thm:x| (for the |amsthm| |proof| environment, for |proof| defined
%   as a theorem in the Springer classes, and for deferred proofs as with
%   \apxproof). When none of this applies, |\proofof| inside the proof pins it to
%   the result labelled \meta{label}. It pins the proof \emph{as a whole},
%   wherever in it it stands, so the references made before it count as much as
%   those made after; \meta{label} may be the label of a result stated later in
%   the document. This is what a proof written in an appendix, or anywhere else
%   away from its result, needs and all that it needs: its dependencies are then
%   captured as those of a proof placed right after its result would be, and
%   recording them by hand with |\uses| or |\proofgraphedge| is not called for.
%
%   The pin covers the proofs \emph{nested} in the pinned one as well, the steps
%   of a structured proof (as \textsf{pf2} writes them) being proofs in their own
%   right; a |\proofof| in a nested proof pins that proof alone, overriding for
%   it the pin of the proof around it. A nested proof that comes after a
%   \emph{result} opened inside the enclosing proof -- the proof of a claim
%   stated in the course of an argument -- attaches to that claim instead, as an
%   ordinary proof would.
% \end{description}
%
% \subsection{Choosing what appears in the graph}
% Result environments are detected automatically; these commands adjust which
% results, and which edges, are kept.
% \begin{description}
% \item[\cmdname{proofgraphtrack}\cmdarg{names}]%
%   \SpecialUsageIndex{\proofgraphtrack}%
%   Forces the listed environment types to be tracked, overriding the default
%   exclusions, which are |definition|, |example|, |remark| and |note|. Use it in
%   the preamble.
% \item[\cmdname{proofgraphuntrack}\cmdarg{names}]%
%   \SpecialUsageIndex{\proofgraphuntrack}%
%   Excludes further environment types, as in |\proofgraphuntrack{observation}|.
%   Use it in the preamble.
% \item[\cmdname{proofgraphproofenv}\cmdarg{names}]%
%   \SpecialUsageIndex{\proofgraphproofenv}%
%   Treats the listed environments as proofs, besides |proof|, so that the
%   references they make become dependencies of the result they establish. Use
%   it in the preamble, as in |\proofgraphproofenv{pf}|. It serves a document
%   whose proofs live in an environment of another name, for instance because a
%   proof package had to be renamed to avoid the clash with \textsf{amsthm}'s
%   |proof| (as with \textsf{pf2}), and equally a |\newtheorem{claimproof}| used
%   for the proofs of claims. Being proofs, the environments named are also
%   removed from the results drawn, even when they are numbered; a
%   |\proofgraphtrack| on the same name overrides that, for an environment to be
%   drawn as a result as well.
%
%   An environment taking an optional title is wrapped as |proof| is, so a
%   detached |\begin{pf}[of Theorem~\ref{thm:x}]| attaches to |thm:x|; one
%   declared with |\newtheorem| attaches to the result it follows, or to the one
%   |\proofof| names.
% \item[\cmdname{proofgraphignore}\cmdarg{A}\cmdarg{B}]%
%   \SpecialUsageIndex{\proofgraphignore}%
%   Drops the dependency of the result \meta{A} on the result \meta{B}, the edge
%   the proof of \meta{A} produced by referring to \meta{B}: an unwanted forward
%   reference, say. It may appear anywhere, the preamble included.
%
%   The two labels are those of |\proofgraphedge|, in that order: the result
%   first, what its proof uses second, the way the dependency is
%   \emph{recorded}. They are \emph{not} the ends of the arrow \emph{drawn},
%   which with the default |direction=usedby| runs the other way, from \meta{B}
%   to \meta{A}. So the rule undoing a |\proofgraphedge{cor:x}{thm:y}| is
%   |\proofgraphignore{cor:x}{thm:y}|, with the same two labels in the same
%   order, and the rule dropping the arrow the picture draws from \texttt{lem:a}
%   to \texttt{thm:m} is |\proofgraphignore{thm:m}{lem:a}|. The |direction|
%   option does not enter into it: like every filter the rule works on the
%   dependency recorded, |direction| settling only how it is drawn, so a graph
%   turned round keeps its rules. A rule that drops nothing is reported at the
%   end of the run, and says so when it is the two labels that have been
%   swapped.
% \item[\cmdname{proofgraphignorecite}\cmdoarg{note}\cmdarg{A}\cmdarg{keys}]%
%   \SpecialUsageIndex{\proofgraphignorecite}%
%   The same for a citation node, which |\proofgraphignore| cannot name, both of
%   its arguments being labels: drops the dependency of the result \meta{A} on
%   the cited works \meta{keys} (a comma-separated list of \BibTeX\ keys). Its
%   arguments are those of |\proofgraphciteedge|, which it undoes. The optional
%   \meta{note} picks, of a work cited with a note and so drawn as more than one
%   node, the node of that note, so
%   |\proofgraphignorecite[Thm~3]{thm:x}{knuth}| leaves the node of a plain
%   |\cite{knuth}| alone. A citation node no edge reaches is not drawn, so
%   dropping the last dependency on a work removes its node with it.
% \item[\cmdname{proofgraphexclude}\cmdarg{label}]%
%   \SpecialUsageIndex{\proofgraphexclude}%
%   Removes the result labelled \meta{label}, together with all its incident
%   edges, from the graph.
% \item[\cmdname{proofgraphcitecommand}\cmdarg{names}]%
%   \SpecialUsageIndex{\proofgraphcitecommand}%
%   Captures further citation commands, given as a comma-separated list of names
%   without their backslash, as in |\proofgraphcitecommand{citeauthor,citeyear}|.
%   Use it in the preamble. Any command taking the usual
%   \meta{cs}|*[|\meta{pre}|][|\meta{post}|]|\marg{keys} arguments works,
%   including one you define yourself. Naming a command the document does not
%   define is harmless.
%
%   It cannot help with the \textsf{biblatex} \emph{multicite} commands
%   (|\cites|, |\parencites|, |\textcites|), whose
%   |(|\meta{pre}|)(|\meta{post}|)| syntax and repeated key groups \proofgraph\
%   does not parse; cite such works individually, or record them with |\uses|.
% \end{description}
%
% \subsection{Styling the graph}
% The attributes are passed verbatim to \Graphviz; see its attribute reference,
% \url{https://graphviz.org/doc/info/attrs.html}, for the node attributes
% available (and \url{https://graphviz.org/doc/info/shapes.html} for the shapes).
% \begin{description}
% \item[\cmdname{proofgraphstyle}\cmdarg{env}\cmdarg{attrs}]%
%   \SpecialUsageIndex{\proofgraphstyle}%
%   Sets \Graphviz\ node attributes for all results of the environment
%   \meta{env}, e.g., |shape=ellipse| or |style=filled,fillcolor=gold|.
% \item[\cmdname{proofgraphstylecite}\cmdarg{attrs}]%
%   \SpecialUsageIndex{\proofgraphstylecite}%
%   Sets the \Graphviz\ attributes of the external citation nodes added by the
%   |cite| option (default |shape=note|), as in
%   |\proofgraphstylecite{shape=box,style=dashed,color=gray}|.
% \end{description}
% A colour may be given in any of the forms \Graphviz\ accepts, including the
% |#|\meta{rrggbb} of an \textsc{rgb} value: the attributes of these two
% commands, and of these two alone, are read with |#| as an ordinary character,
% so |\proofgraphstyle{theorem}{style=filled,fillcolor="#ffd700"}| may be
% written as it stands.
%
% \subsection{Embedding the rendered graph}
% \begin{description}
% \item[\cmdname{proofgraph}\cmdoarg{options}]\SpecialUsageIndex{\proofgraph}%
%   Includes the rendered graph image at this point, passing the optional
%   \meta{options} to the underlying |\includegraphics| (e.g.,
%   |\proofgraph[width=\linewidth]|). This is the only feature that needs
%   \textsf{graphicx}, which the package therefore does not load on its own: a
%   document that uses |\proofgraph| must load \textsf{graphicx} itself (loading
%   \textsf{tikz} for the |hyperlinks| overlay also suffices). With the |autorun|
%   option, \Graphviz\ is run at the end of a compilation in which the graph
%   changed (shell-escape required) to produce that image,
%   so |\proofgraph| embeds the graph from the previous run and is up to date
%   after two runs. Without |autorun|, no \Graphviz\ is run: |\proofgraph| embeds
%   a pre-existing |\jobname-proofgraph.|\meta{format} file if you have rendered
%   one under that name, and otherwise warns. When \textsf{hyperref} and
%   \textsf{TikZ} are both loaded (with |hyperlinks| on and the |-Tplain| file
%   from |autorun| present), every node becomes an internal hyperlink to the
%   result it represents: \Graphviz\ writes the node positions to that |-Tplain|
%   file and |\proofgraph| overlays a transparent |\hyperref| area on each node
%   (needed because the link annotations \Graphviz\ puts in the |.pdf| are
%   discarded by |\includegraphics|).
% \end{description}
%
% \subsection{Options}
% \begin{description}
% \item[\optname{selfloops}]
%   Either |remove| (default), dropping edges from a result to itself, or |keep|,
%   retaining them.
% \item[\optname{autorun}]
%   Boolean (default false). Runs \Graphviz\ automatically at the end of the run;
%   requires compilation with shell-escape enabled. \Graphviz\ is run only when
%   the graph, the engine or the format has changed since the last rendering, so
%   that repeated compilations converge and a build tool such as
%   \textsf{latexmk} settles instead of rerunning indefinitely.
%   Unrestricted shell escape is what is needed, the |-shell-escape| of the
%   usual engines: the restricted mode a \TeX\ installation has by default runs
%   only the few commands the installation lists and \Graphviz\ is not among
%   them. Should the rendering produce nothing, for that reason or because
%   \Graphviz\ is not installed, the package says so at the end of the run.
% \item[\optname{engine}]
%   \Graphviz\ layout engine (default |dot|); any \Graphviz\ engine works, such as
%   |neato|, |fdp|, |sfdp|, |circo| or |twopi| (see
%   \url{https://graphviz.org/docs/layouts/}).
% \item[\optname{format}]
%   Output format for |autorun| and |\proofgraph| (default |pdf|); any \Graphviz\
%   output format works, such as |png|, |svg| or |ps| (see
%   \url{https://graphviz.org/docs/outputs/}).
% \item[\optname{hyperlinks}]
%   Boolean (default true). Makes the nodes of an embedded graph clickable when
%   \textsf{hyperref} and \textsf{TikZ} are available (see |\proofgraph|); set it
%   false to embed the graph as a plain image.
% \item[\optname{direction}]
%   Either |usedby| (default), orienting an edge from the result that is used to
%   the result that uses it (so an arrow |a->b| reads ``|a| is used by |b|''), or
%   |uses|, the reverse.
% \item[\optname{rankdir}]
%   \Graphviz\ graph direction (default |LR|).
% \item[\optname{cite}]
%   Boolean (default false). Also captures |\cite| inside proofs as dependencies
%   on external work, drawn as distinct nodes (Section~\ref{sec:realworld}), and,
%   unless |statements| is off, in the body of a result too. A
%   |\cite| note is kept: |\cite[Thm~3]{key}| labels the node ``Thm~3, \dots''.
%   In the two-optional-argument form of \textsf{natbib} and \textsf{biblatex},
%   |\cite[see][Thm~3]{key}|, the post-note likewise labels the node and the
%   pre-note is ignored. Every citation command of \LaTeX, \textsf{natbib} and
%   \textsf{biblatex} that cites a work is captured the same way, whenever the
%   document defines it, starred variants included: |\citet|, |\citep|,
%   |\citealt|, |\citealp|, |\parencite|, |\textcite|, |\autocite|,
%   |\footcite|, |\footcitetext|, |\smartcite|, |\supercite|, |\fullcite| and
%   |\footfullcite|, together with the commands printing only a fragment of a
%   citation (|\citeauthor|, |\citefullauthor|, |\citetitle|, |\citeyear|,
%   |\citeyearpar|, |\citedate|, |\citeurl|, |\citenum|) and all their
%   capitalised forms. A work invoked twice in one proof, by |\citeauthor| and
%   by |\cite| say, still yields a single edge. |\proofgraphcitecommand| adds
%   any further command; |\usescite| and |\proofgraphciteedge| record a citation
%   node by hand, with or without this option.
% \item[\optname{statements}]
%   Boolean (default true). Captures cross-references, and citations when
%   |cite| is on, in the \emph{body} of a result as well, not only in its proof
%   and in its optional title: a statement reading ``Under the hypotheses of
%   Theorem~2\dots'' or ``\dots\ the bound of \cite[Thm~3]{knuth}'' records
%   that dependency on the result being stated, as a proof making the same
%   reference would. The capture stops at the |\end| of the result, so
%   references in the surrounding text record nothing. Set it false for a graph
%   built from proofs alone, which is the right reading when statements refer to
%   earlier results to contrast with them rather than to rely on them.
% \item[\optname{citelabel}]
%   Either |key| (default), labelling citation nodes by their \BibTeX\ key, or
%   |tag|, using the printed citation tag instead (``[1]'', ``[Knu74]'',
%   ``[Abiteboul et al.\ 1995]''\dots), which is only known after a first run (so
%   it needs two runs). It supports the standard, \textsf{hyperref}-wrapped and
%   \textsf{natbib} (numeric and author--year) bibliographies. \textsf{biblatex}
%   is supported in part: the tag is recovered for its numeric styles (|[1]|) and
%   its alphabetic ones (|[Knu68]|), but not for its author--year styles, whose
%   tag is a formatted name list that \textsf{biblatex} exposes as no single
%   field. Whenever no clean tag can be recovered the node keeps the key, which
%   is always a correct, if terser, label.
%
%   The note of a |\cite| is set where a citation prints it, inside the tag and
%   after it: |\cite[pp.~23--27]{knuth}| gives a node labelled ``[Knu68,
%   pp.~23--27]'', carrying the dashes of the document, and ``knuth, pp.~23--27''
%   when the label is the key.
% \item[\optname{file}]
%   Base name of the output file (default |\jobname|).
% \end{description}
%
% \subsection{Use with \apxproof}
% \label{sec:apxproof}
% \apxproof\ moves proofs to the appendix and typesets them there, so the
% references a proof makes are only executed at that later point.
% \textsf{proofgraph} therefore captures each proof where \apxproof\ actually
% typesets it: a proof deferred by a repeated-theorem environment is attributed
% to its theorem through the hook |\apxproofhook|, which \apxproof\ fires
% inside each such proof, and a plain |proof| -- left in the body, or belonging
% to a |toappendix| block and replayed with it -- to the result it follows, as
% without \apxproof. The hook is available in \apxproof\ \textbf{v1.4.1} and
% later; load the two packages in either order and no further setup is needed.
%
% With an \emph{earlier} \apxproof\ the document still compiles cleanly and
% every result still becomes a node, but \emph{no proof dependencies are
% captured}, because \textsf{proofgraph} cannot tell which appendix proof
% belongs to which result. In that case, add them manually with |\uses| (and
% |\proofof| to pin a proof to its result), or upgrade \apxproof.
%
% \section{Limitations}
% \begin{itemize}
% \item Environments declared through |\newtheorem| (or |\spnewtheorem|)
%   \emph{after} the package is loaded are tracked automatically. Environments
%   predefined by the document class \emph{before} that (as in \textsf{lipics},
%   \textsf{acmart} or the Springer classes) are also detected, provided their
%   name is one of the common result names \proofgraph\ probes at
%   |\begin{document}| (|theorem|, |lemma|, |proposition|, |corollary|\dots); an
%   environment with an unusual name predefined before the package is loaded must
%   be added with |\proofgraphtrack|.
% \item Results must be labelled for references to them to be resolved.
% \item Two compilation runs are required, as for any cross-reference.
% \end{itemize}
%
% \section{See also}
% \label{sec:seealso}
% \textsf{result-graph}~\cite{resultgraph} does for a \textsf{Lean}~4
% formalisation what \proofgraph\ does for an article: it draws the dependency
% graph of the results of a development, deliberately with the same vocabulary,
% node styles and edge convention, so that the graph of a formalisation and the
% graph of the paper it formalises look like one family. There the dependencies
% are not inferred from cross-references but read off the elaborated proof terms,
% and are therefore exact; they also distinguish, for free, a result used in the
% \emph{statement} of another from one used only in its \emph{proof}, which
% \proofgraph, under the |statements| option, records but draws alike. And since
% a whole development is far too
% large for one figure, it is sliced into a small neighbourhood of each result,
% rather than drawn as a single graph as here.
%
% \begin{thebibliography}{9}
% \bibitem{euclid} Euclid. \emph{Elements}, Book~IX. c.~300~BCE.
% \bibitem{AmarilliCMS20} A.~Amarilli, F.~Capelli, M.~Monet, and P.~Senellart.
%   Connecting knowledge compilation classes and width parameters.
%   \emph{Theory of Computing Systems}, \textbf{64}(5):861--914, 2020.
%   \href{https://doi.org/10.1007/s00224-019-09930-2}%
%   {\texttt{doi:10.1007/s00224-019-09930-2}}.
% \bibitem{resultgraph} P.~Senellart. \emph{result-graph: dependency graphs of
%   the results of a Lean~4 project}.
%   \url{https://github.com/PierreSenellart/result-graph}.
% \end{thebibliography}
%
% \StopEventually{
%   \clearpage
%   \PrintIndex
%   \PrintChanges
% }
%
% \clearpage
% \section{Implementation}
%
% The format required is that of October 2020, the first to provide
% |\NewDocumentCommand| in the kernel; \textsf{expl3}, which the package also
% uses, has been preloaded since the February 2020 release.
%    \begin{macrocode}
%<*package>
\NeedsTeXFormat{LaTeX2e}[2020/10/01]
\ProvidesPackage{proofgraph}
  [2026/07/31 v1.1.0 Dependency graph of results]
%    \end{macrocode}
%
% \subsection{Dependencies}
% \textsf{amsthm} (loaded for its heading hooks) defines the |proof| environment
% with |\newenvironment|, which aborts with ``Command |\proof| already defined''
% under a class that already provides one (e.g.\ the Springer |svjour3| |\proof|).
% When we are about to load \textsf{amsthm} ourselves (it is not already loaded),
% we save the class's |proof|/|\endproof| and free them so \textsf{amsthm} loads
% cleanly, then restore the class's versions afterwards, so its proof styling is
% kept (and \proofgraph\ still hooks that |proof| environment). If \textsf{amsthm}
% is already loaded, |\RequirePackage| is a no-op and |\proof| is left untouched.
%    \begin{macrocode}
\@ifpackageloaded{amsthm}{}{%
  \@ifundefined{proof}{}{%
    \let\pgph@class@proof\proof \let\pgph@class@endproof\endproof}%
  \let\proof\relax \let\endproof\relax}
\RequirePackage{amsthm}
\@ifundefined{pgph@class@proof}{}{%
  \let\proof\pgph@class@proof \let\endproof\pgph@class@endproof}
\RequirePackage{etoolbox}
\RequirePackage{xstring}
\RequirePackage{kvoptions}
\RequirePackage{pdftexcmds}
\RequirePackage{shellesc}
%    \end{macrocode}
% \textsf{pdftexcmds} provides |\pdf@filemdfivesum|, with which |autorun|
% recognises an unchanged graph and skips the \Graphviz\ run
% (see |\pgph@maybrender|). On an engine that offers no file checksum the
% package still works, only without that optimisation.
%
% \textsf{shellesc} provides |\ShellEscape|, with which |autorun| runs
% \Graphviz, and |\ShellEscapeStatus|, which says whether it may. The
% primitive |\immediate\write18| is not a portable way to run a command: under
% \LuaTeX\ stream~18 is an ordinary output stream like any other, so the text
% of the command is written out instead of being executed, silently and
% whatever |-shell-escape| was passed. |\ShellEscape| is |\immediate\write18|
% where that means what it says and |os.execute| under \LuaTeX. This is why a
% document that loaded \textsf{minted}, which loads \textsf{shellesc} too,
% used to render its graph under \LuaLaTeX\ while the same document without it
% did not.
%
% \textsf{graphicx} is needed only to embed a rendered graph with |\proofgraph|;
% the core |.dot|-writing machinery does not use it. Rather than burden every
% document with it, we leave it unloaded and check for |\includegraphics| at the
% point of use (see |\proofgraph|). When the |hyperlinks| overlay is active
% \textsf{tikz} is loaded, which pulls \textsf{graphicx} in anyway.
% \begin{macro}{\pgph@trimsp}
% An expandable space-trimmer (labels inside a |\cite|/|\cref|/|\uses| list
% arrive with the space after each comma), built on the |expl3| primitive so it
% is robust to leading spaces and braces.
%    \begin{macrocode}
\ExplSyntaxOn
\cs_new:Npn \pgph@trimsp #1 { \tl_trim_spaces:n {#1} }
\ExplSyntaxOff
%    \end{macrocode}
% \end{macro}
%
% \subsection{Options}
%    \begin{macrocode}
\SetupKeyvalOptions{family=pgph,prefix=pgph@}
\DeclareStringOption[remove]{selfloops}
\DeclareBoolOption{autorun}
\DeclareStringOption[dot]{engine}
\DeclareStringOption[pdf]{format}
\DeclareStringOption[LR]{rankdir}
\DeclareStringOption[usedby]{direction}
\DeclareBoolOption{cite}
\DeclareStringOption[key]{citelabel}
\DeclareBoolOption[true]{statements}
\DeclareBoolOption[true]{hyperlinks}
\DeclareStringOption{file}
\ProcessKeyvalOptions*
\ifx\pgph@file\@empty\def\pgph@file{\jobname}\fi
%    \end{macrocode}
%
% \subsection{Literal braces}
% We need catcode-12 braces to write the \Graphviz\ |digraph { }| block.
%    \begin{macrocode}
\begingroup
  \catcode`\[=1 \catcode`\]=2
  \catcode`\{=12 \catcode`\}=12
  \gdef\pgph@ob[{]%
  \gdef\pgph@cb[}]%
\endgroup
%    \end{macrocode}
%
% \subsection{Data structures}
% A global counter gives each result instance a unique numeric id. Nodes
% and edges are accumulated in list macros and emitted at end of
% document, so that editing commands (ignore, exclude, self-loops) can
% be applied to the complete graph regardless of where they appear.
%    \begin{macrocode}
\newcount\pgph@idcnt
\newcount\pgph@proofcnt
\newcount\pgph@citeslotcnt
\def\pgph@nodes{}
\def\pgph@edges{}
\def\pgph@ignores{}
\def\pgph@excludes{}
\let\pgph@curresult\@empty
\let\pgph@curproof\@empty
%    \end{macrocode}
%
% \begin{macro}{\pgph@auxmap}
% The label-to-id map. |\label| inside a result records, both in memory
% and in the |.aux| file, that the user label maps to the current node
% id, so references resolve independently of \textsf{hyperref}.
%    \begin{macrocode}
\newcommand\pgph@auxmap[2]{\global\@namedef{pgph@map@#1}{#2}%
  \ifcsname pgph@rmap@#2\endcsname\else
    \global\@namedef{pgph@rmap@#2}{#1}\fi}
\newcommand\pgph@setmap[2]{%
  \edef\pgph@tmpid{#2}%
  \global\expandafter\edef\csname pgph@map@#1\endcsname{\pgph@tmpid}%
  \ifcsname pgph@rmap@\pgph@tmpid\endcsname\else
    \global\expandafter\def\csname pgph@rmap@\pgph@tmpid\endcsname{#1}%
  \fi
  \protected@write\@auxout{}{\string\pgph@auxmap{#1}{\pgph@tmpid}}%
}
%    \end{macrocode}
% \end{macro}
%
% \subsection{Registering result environments}
% \begin{macro}{\pgph@register}
% A result environment is tracked by installing a begin-hook that records a
% node. In every declaration command (|\newtheorem|, |\spnewtheorem|) the
% environment \emph{name} is the first mandatory argument, so registration
% does not need to parse the remaining (varied) arguments. Each name met is
% queued as a \emph{candidate} (deduplicated); the begin-hooks themselves are
% installed only at |\begin{document}|, once the inclusion and exclusion lists
% are settled, so the customisation commands below may appear anywhere in the
% preamble.
%    \begin{macrocode}
\let\pgph@candidates\@empty
\newcommand\pgph@register[1]{%
  \edef\pgph@thisname{#1}%
  \IfBeginWith\pgph@thisname{axp@}{}{%
    \ifcsname pgph@cand@\pgph@thisname\endcsname\else
      \global\@namedef{pgph@cand@\pgph@thisname}{}%
      \ifx\pgph@candidates\@empty
        \global\let\pgph@candidates\pgph@thisname
      \else
        \xdef\pgph@candidates{\pgph@candidates,\pgph@thisname}%
      \fi
    \fi
  }%
}
%    \end{macrocode}
% \end{macro}
%
% We wrap |\newtheorem| (and, for Springer classes, |\spnewtheorem|) to
% register each newly declared environment, replaying the original command so
% its own argument parsing is untouched. The starred (unnumbered) forms are
% \emph{not} auto-registered: a results graph wants numbered statements, and
% these forms are commonly used for unnumbered repeats of a result (e.g., the
% appendix copies produced by \apxproof\ and \textsf{myproofs}), which would
% otherwise appear as spurious unnumbered duplicate nodes. Use |\proofgraphtrack|
% to track an unnumbered environment deliberately.
%    \begin{macrocode}
\let\pgph@orig@newtheorem\newtheorem
\renewcommand\newtheorem{\@ifstar{\pgph@nt@star}{\pgph@nt@plain}}
\newcommand\pgph@nt@plain[1]{%
  \pgph@register{#1}\pgph@orig@newtheorem{#1}}
\newcommand\pgph@nt@star[1]{\pgph@orig@newtheorem*{#1}}
\ifdefined\spnewtheorem
  \let\pgph@orig@spnewtheorem\spnewtheorem
  \renewcommand\spnewtheorem{%
    \@ifstar{\pgph@spnt@star}{\pgph@spnt@plain}}
  \newcommand\pgph@spnt@plain[1]{%
    \pgph@register{#1}\pgph@orig@spnewtheorem{#1}}
  \newcommand\pgph@spnt@star[1]{\pgph@orig@spnewtheorem*{#1}}
\fi
%    \end{macrocode}
%
% \DescribeMacro{\proofgraphtrack}
% \DescribeMacro{\proofgraphuntrack}
% \begin{macro}{\proofgraphtrack}
% \begin{macro}{\proofgraphuntrack}
% Result environments are detected automatically, but the analysis can be
% tuned. |\proofgraphtrack{|\meta{names}|}| forces the listed environments to be
% tracked, overriding any exclusion; |\proofgraphuntrack{|\meta{names}|}|
% excludes the listed environment \emph{types}. Expository environments are
% excluded by default: |definition|, |example|, |remark|, |note|. Both commands
% belong in the preamble. |\pgph@defaultenvs| is the list of class-predefined
% names probed at |\begin{document}|; |\pgph@defaultexclude| seeds the exclusion
% set.
%    \begin{macrocode}
\newcommand\pgph@defaultenvs{theorem,lemma,proposition,corollary,%
  definition,conjecture,claim,fact,remark,example,observation,%
  notation,problem,question,property,assumption,hypothesis,principle}
\newcommand\pgph@defaultexclude{definition,example,remark,note}
\newcommand\proofgraphtrack[1]{%
  \@for\pgph@one:=#1\do{%
    \edef\pgph@one{\expandafter\pgph@trimsp\expandafter{\pgph@one}}%
    \csundef{pgph@xtype@\pgph@one}%
    \global\@namedef{pgph@force@\pgph@one}{}%
    \expandafter\pgph@register\expandafter{\pgph@one}%
  }%
}
\newcommand\proofgraphuntrack[1]{%
  \@for\pgph@one:=#1\do{%
    \edef\pgph@one{\expandafter\pgph@trimsp\expandafter{\pgph@one}}%
    \csundef{pgph@force@\pgph@one}%
    \global\@namedef{pgph@xtype@\pgph@one}{}%
  }%
}
\expandafter\proofgraphuntrack\expandafter{\pgph@defaultexclude}
%    \end{macrocode}
% \end{macro}
% \end{macro}
%
% Environments predefined by the document class (\textsf{acmart},
% \textsf{lipics}, Springer classes) are declared before the package is loaded,
% so the wrappers above never see them. At |\begin{document}| we add every
% common result name that exists to the candidate list, then install a
% begin-hook for each candidate that is forced or not excluded.
%    \begin{macrocode}
\newcommand\pgph@dohook[1]{%
  \ifcsname pgph@hooked@#1\endcsname\else
    \global\@namedef{pgph@hooked@#1}{}%
    \AtBeginEnvironment{#1}{\pgph@beginresult{#1}}%
  \fi
}
\newcommand\pgph@maybehook[1]{%
  \ifcsname pgph@force@#1\endcsname
    \pgph@dohook{#1}%
  \else
    \ifcsname pgph@xtype@#1\endcsname\else\pgph@dohook{#1}\fi
  \fi
}
\newcommand\pgph@resolveregistration[1]{%
  \@for\pgph@one:=#1\do{%
    \ifcsname\pgph@one\endcsname%
      \expandafter\pgph@register\expandafter{\pgph@one}\fi}%
  \ifx\pgph@candidates\@empty\else
    \@for\pgph@one:=\pgph@candidates\do{%
      \expandafter\pgph@maybehook\expandafter{\pgph@one}}%
  \fi
}
\AtBeginDocument{%
  \expandafter\pgph@resolveregistration\expandafter{\pgph@defaultenvs}}
%    \end{macrocode}
%
% \begin{macro}{\pgph@beginresult}
% \begin{macro}{\pgph@hooknames}
% \begin{macro}{\pgph@recordresult}
% Start a result: allocate an id and record the node. The node is keyed and
% styled by the environment name (third argument, also a fallback label), but
% displayed by the result's \emph{formatted} name (``Theorem'', ``Lemma''),
% captured from the heading. We get that name by locally wrapping the heading
% macros, whose first two arguments are the formatted name and the number. Which
% macro carries the \emph{optional note} (where a restatement's cross-reference
% lives) depends on the setup: \textsf{amsthm}, which \proofgraph\ always loads,
% routes every heading through |\@begintheorem|, defined as |#1#2[#3]| with the
% note as the bracketed |#3| (and leaves |\@opargbegintheorem| |\relax|); the
% \LaTeX\ kernel instead uses a two-argument |\@begintheorem| with a separate
% three-argument |\@opargbegintheorem| for the note; the Springer classes use
% |\@spbegintheorem|/|\@spopargbegintheorem|. We detect the \textsf{amsthm} form
% (|\@opargbegintheorem| is |\relax|) and read the bracketed note there,
% otherwise we hook |\@opargbegintheorem|; either way the note goes to
% |\pgph@scanrefs|. These redefinitions are confined to the current environment
% group, so they revert at |\end|. We also locally redefine |\label| to populate
% the map.
%
% The note read as a \emph{delimited} argument is handed back to the original
% macro inside braces, |[{|\meta{note}|}]| rather than |[|\meta{note}|]|. TeX
% strips the outer braces of a delimited argument that is one brace group in
% its entirety, so a title written |[{\cite[p.~37]{a-l}}]| -- the bracing the
% \textsf{amsthm} manual prescribes to hide the |]| of the citation note -- would
% otherwise reach the original macro unbraced, whose own scan for |]| would then
% stop inside the citation and take the rest of the note for text. Rebracing
% restores exactly what was written: the added braces are stripped again by that
% scan.
%
% Nodes are accumulated in \emph{reverse} order of appearance
% (|\pgph@prependnode|). With |rankdir=LR|, \Graphviz\ stacks disconnected
% components vertically in the order they are read, from the bottom up; emitting
% the results last-to-first therefore lays them out top-to-bottom in document
% order (first result on top), while connected results stay grouped. This is a
% layout tendency, not a guarantee, but needs no |pack| trickery.
%
% Unless |statements| is switched off, the \emph{body} of the result is treated
% like the body of a proof of that result (|\pgph@beginstatement|), so that a
% cross-reference or a citation made in the statement itself becomes a
% dependency. Those assignments are local to the environment group, as the ones
% above are, so they revert at |\end|.
%
% A result opened \emph{inside} a proof -- a claim stated in the course of an
% argument, and proved there -- is a step of that proof, so |\pgph@supportedge|
% records the proof as depending on it: the claim would otherwise be drawn with
% whatever its own proof uses hanging below it and nothing above, as a component
% detached from the result it was stated to establish. The edge is recorded
% against the enclosing proof rather than a result, so it follows the same chain
% of owners as any other, |\proofof| and nesting included, and it passes through
% the ordinary filters, so |\proofgraphignore| drops it like any other. Nothing
% is recorded for a claim the document leaves unnumbered: |\pgph@node| excludes
% such a result, and the edge to it goes with it.
%    \begin{macrocode}
\newcommand\pgph@supportedge{%
  \ifx\pgph@encproof\@empty\else
    \pgph@addedge{\pgph@encproof}{node:\pgph@curresult}%
  \fi
}
\newcommand\pgph@prependnode[1]{%
  \xdef\pgph@nodes{\unexpanded{#1}\unexpanded\expandafter{\pgph@nodes}}%
}
\newcommand\pgph@beginresult[1]{%
  \let\pgph@encproof\pgph@curproof
  \global\advance\pgph@idcnt\@ne
  \xdef\pgph@curresult{\the\pgph@idcnt}%
  \protected@edef\pgph@tmp{%
    \noexpand\pgph@node{\pgph@curresult}{#1}{#1}}%
  \expandafter\pgph@prependnode\expandafter{\pgph@tmp}%
  \pgph@supportedge
  \pgph@hooknames
  \ifx\label\pgph@maplabel\else\let\pgph@savedlabel\label\fi
  \let\label\pgph@maplabel
  \ifpgph@statements\pgph@beginstatement\fi
}
\newcommand\pgph@hooknames{%
  \ifdefined\@begintheorem
    \let\pgph@orig@bt\@begintheorem
    \ifx\@opargbegintheorem\relax
      \def\@begintheorem##1##2[##3]{%
        \pgph@recordresult{##1}{##2}\pgph@scanrefs{##3}%
        \pgph@orig@bt{##1}{##2}[{##3}]}%
    \else
      \def\@begintheorem##1##2{%
        \pgph@recordresult{##1}{##2}\pgph@orig@bt{##1}{##2}}%
    \fi
  \fi
  \ifdefined\@opargbegintheorem\ifx\@opargbegintheorem\relax\else
    \let\pgph@orig@obt\@opargbegintheorem
    \def\@opargbegintheorem##1##2##3{%
      \pgph@recordresult{##1}{##2}\pgph@scanrefs{##3}%
      \pgph@orig@obt{##1}{##2}{##3}}%
  \fi\fi
  \ifdefined\@spbegintheorem
    \let\pgph@orig@spbt\@spbegintheorem
    \def\@spbegintheorem##1##2##3##4{%
      \pgph@recordresult{##1}{##2}\pgph@orig@spbt{##1}{##2}{##3}{##4}}%
  \fi
  \ifdefined\@spopargbegintheorem
    \let\pgph@orig@spobt\@spopargbegintheorem
    \def\@spopargbegintheorem##1##2##3##4##5{%
      \pgph@recordresult{##1}{##2}\pgph@scanrefs{##3}%
      \pgph@orig@spobt{##1}{##2}{##3}{##4}{##5}}%
  \fi
}
%    \end{macrocode}
% \begin{macro}{\pgph@scanrefs}
% A result whose optional title carries a cross-reference, e.g., a |result|
% whose title is |[(Corollary~\ref{cor:x})]|, is understood to
% restate or build on that result: we record an edge from the current result to
% each label referenced in the title. We capture these by typesetting the title
% once into a discarded box with the capturing commands installed and
% |\pgph@curproof| pointing at the current result. Everything happens inside a
% group and |\pgph@curproof| is only \emph{read} while the box is built, so a
% local assignment suffices and nothing outside the title is affected; a global
% one would escape the enclosing environment, which under |statements| (where
% |\pgph@curproof| is already set on entering the result) would leave every
% subsequent reference attributed to that result.
%    \begin{macrocode}
\newcommand\pgph@scanrefs[1]{%
  \ifx\pgph@curresult\@empty\else
    \begingroup
      \let\pgph@curproof\pgph@curresult
      \pgph@installcapture
      \setbox\z@\hbox{#1}%
    \endgroup
  \fi
}
%    \end{macrocode}
% \end{macro}
% The formatted name is the first argument of every heading macro
% (|\@begintheorem|/|\@opargbegintheorem| for |amsthm|/|\newtheorem|,
% |\@spbegintheorem|/|\@spopargbegintheorem| for Springer's |\spnewtheorem|);
% the second argument is the number. We use that second argument only to tell
% \emph{whether} the result is numbered (it is empty for the starred,
% unnumbered forms): an unnumbered result records no number, so it is dropped at
% emission unless tracked explicitly. The number \emph{value}, when present, is
% read from |\@currentlabel| rather than from that argument, which some classes
% pollute (e.g., \textsf{lipics} injects |\advance\par@deathcycles\@ne|).
% |\@currentlabel|, set by the result's |\refstepcounter| just before the
% heading, is the clean number and needs no |\label|. Values are reduced to plain
% strings now, with fragile wrappers neutralized (notably \apxproof's
% forward-link). The first heading wins (guarded on the name), ignoring nested
% headings. The emptiness of the number argument is tested without expanding it,
% so a polluted (but non-empty) number does not trip the test.
%    \begin{macrocode}
\newcommand\pgph@recordresult[2]{%
  \ifcsname pgph@name@\pgph@curresult\endcsname\else
    \begingroup
      \let\protect\@empty
      \def\axp@forward@link##1##2{##2}%
      \edef\pgph@tmpname{#1}%
      \global\expandafter\let\csname pgph@name@\pgph@curresult\endcsname
                              \pgph@tmpname
      \global\@namedef{pgph@seen@\pgph@curresult}{}%
      \def\pgph@tmphn{#2}%
      \ifx\pgph@tmphn\@empty\else
        \edef\pgph@tmpnum{\@currentlabel}%
      \global\expandafter\let\csname pgph@num@\pgph@curresult%
        \endcsname\pgph@tmpnum
      \fi
    \endgroup
  \fi
}
%    \end{macrocode}
% \end{macro}
% \end{macro}
% \end{macro}
% |\pgph@maplabel| records the label-to-node mapping. As a fallback for classes
% whose headings pass through none of the hooked heading macros (so no heading
% was |seen|), it also captures the number from |\@currentlabel|. We do
% \emph{not} do this once a heading was seen: there, an absent number means the
% result is genuinely unnumbered (a starred form), and |\@currentlabel| could be
% a stale value from an earlier result. The value is expanded to a plain string
% \emph{now}, inside a group where fragile wrappers are neutralized, so that
% nothing dangerous survives to the emission pass: in particular, under
% \textsf{apxproof} forward-linking, |\@currentlabel| holds
% |\axp@forward@link{|\meta{target}|}{|\meta{number}|}|, whose hyperlink machinery
% must never be expanded in a non-typesetting context (it would run away).
% Reducing it to its number argument keeps it harmless.
% \begin{macro}{\pgph@maplabel}
%    \begin{macrocode}
\newcommand\pgph@maplabel[1]{%
  \pgph@setmap{#1}{\pgph@curresult}%
  \ifcsname pgph@num@\pgph@curresult\endcsname\else
   \ifcsname pgph@seen@\pgph@curresult\endcsname\else
    \begingroup
      \let\protect\@empty
      \def\axp@forward@link##1##2{##2}%
      \edef\pgph@tmpnum{\@currentlabel}%
      \global\expandafter\let\csname pgph@num@\pgph@curresult%
        \endcsname\pgph@tmpnum
    \endgroup
   \fi
  \fi
  \pgph@savedlabel{#1}%
}
%    \end{macrocode}
% \end{macro}
%
% \subsection{Hooking proofs}
% Inside a proof we redefine the reference commands to capture edges from
% the proved result.
%    \begin{macrocode}
\newcommand\pgph@addedge[2]{%
  \protected@edef\pgph@tmpedge{\noexpand\pgph@edge{#1}{#2}}%
  \expandafter\g@addto@macro\expandafter\pgph@edges%
    \expandafter{\pgph@tmpedge}%
}
\newcommand\pgph@recordref[1]{%
  \ifpgph@nocapture\else
  \ifx\pgph@curproof\@empty\else
    \@for\pgph@one:=#1\do{%
      \edef\pgph@k{\expandafter\pgph@trimsp\expandafter{\pgph@one}}%
      \pgph@addedge{\pgph@curproof}{\pgph@k}}%
  \fi\fi
}
%    \end{macrocode}
% \begin{macro}{\ifpgph@nocapture}
% The capture is installed for the whole of a proof, and a proof long enough to
% break the page has the page furniture built while it is open. The running head
% is set from the marks, so a cross-reference or a citation in a sectional title
% is \emph{executed again} there, in the output routine, and would be recorded
% as a dependency of the proof the break fell in: a document whose section title
% names a result gained an edge from that result to whichever result the proof
% belonged to. Recording is therefore suspended while the page is assembled.
%
% The kernel neutralizes |\label|, |\index| and |\glossary| at the same point,
% in |\@outputpage|, and offers the hook |build/page/reset| there, inside the
% group that encloses the shipped box, so a local switch set from the hook
% covers the head and foot and is undone by itself. Where the hook is not
% declared (a kernel older than the hook, a class with an output routine of its
% own that skips it), |\@outputpage| is wrapped instead, which needs the switch
% to be set globally as the group ends inside the box being shipped.
%    \begin{macrocode}
\newif\ifpgph@nocapture
\newcommand\pgph@wrapoutputpage{%
  \let\pgph@orig@outputpage\@outputpage
  \def\@outputpage{%
    \global\pgph@nocapturetrue
    \pgph@orig@outputpage
    \global\pgph@nocapturefalse}%
}
\AtBeginDocument{%
  \ifdefined\IfHookExistsTF
    \IfHookExistsTF{build/page/reset}%
      {\AddToHook{build/page/reset}{\pgph@nocapturetrue}}%
      {\pgph@wrapoutputpage}%
  \else
    \pgph@wrapoutputpage
  \fi
}
%    \end{macrocode}
% \end{macro}
%
% \begin{macro}{\pgph@beginproof}
% Attach the proof to the current result and install the capturing
% versions of the reference commands (only those that exist).
%
% The proof is not recorded as the result it proves but as a \emph{proof
% instance} of its own, |p|\meta{n}, whose owning result is held in
% |\pgph@powner@p|\meta{n} and read only at emission. An edge captured in the
% proof therefore does not fix its source when it is seen, which is what lets
% |\proofof| pin the whole proof and not merely the references that follow it.
% Were the source fixed at capture, a proof opening with ``By Theorem~1'' and
% naming its own result only at the end would give that first reference to
% whichever result happens to precede the proof. The indirection also lets a
% proof be pinned to a result labelled later in the document, the label being
% resolved with all the others when the graph is written.
%    \begin{macrocode}
%    \end{macrocode}
% The capturing commands are |\protected| so that, when a reference appears
% in a moving argument inside a proof (e.g., a |\caption|), they do not
% expand into the |.aux|/|.lof| and leak internal tokens.
%    \begin{macrocode}
\protected\def\pgph@cap@ref#1{\pgph@recordref{#1}\pgph@saved@ref{#1}}
\protected\def\pgph@cap@cref#1{\pgph@recordref{#1}\pgph@saved@cref{#1}}
\protected\def\pgph@cap@Cref#1{\pgph@recordref{#1}\pgph@saved@Cref{#1}}
\protected\def\pgph@cap@autoref#1{%
  \pgph@recordref{#1}\pgph@saved@autoref{#1}}
\protected\def\pgph@cap@eqref#1{%
  \pgph@recordref{#1}\pgph@saved@eqref{#1}}
\protected\def\pgph@cap@vref#1{\pgph@recordref{#1}\pgph@saved@vref{#1}}
%    \end{macrocode}
% \textsf{zref-clever}'s |\zcref| is the one captured command that is not simply
% \meta{cs}\marg{labels}: it is |\zcref*[|\meta{options}|]|\marg{labels}, and
% both the star and the options have to be handed on unchanged. Its mandatory
% argument is a comma-separated list, which |\pgph@recordref| already splits.
%    \begin{macrocode}
\NewDocumentCommand\pgph@cap@zcref{sO{}m}{%
  \pgph@recordref{#3}%
  \IfBooleanTF{#1}{\pgph@saved@zcref*[#2]{#3}}{\pgph@saved@zcref[#2]{#3}}%
}
\newcommand\pgph@install[1]{%
  \ifcsname #1\endcsname
    \expandafter\ifx\csname #1%
      \expandafter\endcsname\csname pgph@cap@#1\endcsname
    \else
      \expandafter\let\csname pgph@saved@#1%
        \expandafter\endcsname\csname #1\endcsname
      \expandafter\let\csname #1%
        \expandafter\endcsname\csname pgph@cap@#1\endcsname
    \fi
  \fi
}
%    \end{macrocode}
% The set of captured commands is needed in three places (a proof, the optional
% title of a result, and an \apxproof\ deferred proof), so it is named once
% here: adding a command to the list must not be a matter of remembering three
% call sites.
%    \begin{macrocode}
\newcommand\pgph@installrefs{%
  \pgph@install{ref}\pgph@install{cref}\pgph@install{Cref}%
  \pgph@install{autoref}\pgph@install{eqref}\pgph@install{vref}%
  \pgph@install{zcref}%
}
%    \end{macrocode}
% The cross-reference commands and, when the |cite| option is on, the citation
% commands are always installed together, so the pair has a name of its own:
% every place that opens a capturing context (a proof, the optional title of a
% result, and, with |statements|, the body of a result) installs exactly the
% same set.
%    \begin{macrocode}
\newcommand\pgph@installcapture{%
  \pgph@installrefs
  \ifpgph@cite\pgph@installcite\fi
}
%    \end{macrocode}
% A proof may sit inside another one, either because the proof of a claim stated
% in the course of a proof is itself a |proof| environment, or because the
% document writes structured proofs, whose steps are proofs nested in the proof
% of the whole (as \textsf{pf2} does). The two cases want different owners, and
% what tells them apart is whether a \emph{result} has been opened since the
% enclosing context began: the proof of a claim follows the claim, whereas the
% step of a structured proof follows nothing. So each capturing context records,
% in |\pgph@pres@|\meta{ctx}, the result current when it began; a proof opening
% while that result is still the current one is a \emph{sub-proof} and takes the
% enclosing context as its owner, and any other proof takes the current result,
% as before.
%
% Owning the enclosing context rather than a result is what lets a |\proofof| in
% an outer proof reach the inner ones: the pin is read where the chain of owners
% is followed, at emission, so it applies to a nested reference just as it does
% to one made directly in the pinned proof. An inner |\proofof| still pins its
% own proof alone, which is what the detached proof of a claim needs.
%    \begin{macrocode}
\newcommand\pgph@setpres[1]{%
  \global\expandafter\let\csname pgph@pres@#1\endcsname\pgph@curresult
}
\newcommand\pgph@setpowner[2]{%
  \global\expandafter\let\csname pgph@powner@#1\endcsname#2%
}
\newcommand\pgph@beginproof{%
  \global\advance\pgph@proofcnt\@ne
  \edef\pgph@newproof{p\the\pgph@proofcnt}%
  \pgph@setpres\pgph@newproof
  \edef\pgph@ncur{\pgph@curresult}%
  \edef\pgph@npres{%
    \ifcsname pgph@pres@\pgph@curproof\endcsname
      \csname pgph@pres@\pgph@curproof\endcsname
    \else\pgph@nomatch\fi}%
  \ifx\pgph@ncur\pgph@npres
    \pgph@setpowner\pgph@newproof\pgph@curproof
  \else
    \pgph@setpowner\pgph@newproof\pgph@curresult
  \fi
  \let\pgph@curproof\pgph@newproof
  \pgph@installcapture
  \pgph@hookproofhead
}
%    \end{macrocode}
% The sentinel stands for ``no enclosing context'', and must be something no
% result id can be, so that the comparison fails rather than matching an empty
% |\pgph@curresult| outside any result.
%    \begin{macrocode}
\newcommand\pgph@nomatch{-}
%    \end{macrocode}
% \begin{macro}{\pgph@beginstatement}
% Under the |statements| option the body of a result captures like the body of
% its proof. The proof-heading hook is deliberately \emph{not} installed here: it
% turns a cross-reference in an optional title into a |\proofof|, which is right
% for a proof titled ``of Theorem~1'' but not for a result titled ``(restates
% Theorem~1)'', where the reference is an ordinary dependency.
%    \begin{macrocode}
\newcommand\pgph@beginstatement{%
  \let\pgph@curproof\pgph@curresult
  \pgph@setpres\pgph@curproof
  \pgph@installcapture
}
%    \end{macrocode}
% \end{macro}
% When the |proof| environment is itself a theorem-like environment (as with
% Springer's |\spnewtheorem*{proof}|), its optional title is set by
% |\@opargbegintheorem|/|\@spopargbegintheorem|. We wrap these for the duration
% of the proof so that a title such as |[of Theorem~\ref{thm:x}]| sets
% |\proofof{thm:x}|, attributing the proof to its result. As in
% |\pgph@hooknames|, a |\relax| |\@opargbegintheorem| is left alone: that is the
% \textsf{amsthm} setup, where the macro is never called, and giving it a
% meaning would make |\pgph@hooknames| read a result nested in the proof as
% having the kernel's two-argument heading rather than \textsf{amsthm}'s.
%    \begin{macrocode}
\newcommand\pgph@hookproofhead{%
  \ifdefined\@opargbegintheorem\ifx\@opargbegintheorem\relax\else
    \let\pgph@orig@obtp\@opargbegintheorem
    \def\@opargbegintheorem##1##2##3{%
      \pgph@scanproofarg{##3}\pgph@orig@obtp{##1}{##2}{##3}}%
  \fi\fi
  \ifdefined\@spopargbegintheorem
    \let\pgph@orig@spobtp\@spopargbegintheorem
    \def\@spopargbegintheorem##1##2##3##4##5{%
      \pgph@scanproofarg{##3}%
      \pgph@orig@spobtp{##1}{##2}{##3}{##4}{##5}}%
  \fi
  \ifdefined\@Opargbegintheorem
    \let\pgph@orig@Obtp\@Opargbegintheorem
    \def\@Opargbegintheorem##1##2##3##4{%
      \pgph@scanproofarg{##2}\pgph@orig@Obtp{##1}{##2}{##3}{##4}}%
  \fi
}
%    \end{macrocode}
% \end{macro}
% \begin{macro}{\pgph@scanproofarg}
% \begin{macro}{\pgph@wrapproof}
% A proof whose optional title names its result (a detached proof) has its
% result given by that title, e.g., \texttt{thm:x} for an optional argument
% |[of Theorem~\ref{thm:x}]|, rather than by the most recently opened result.
% We capture this by wrapping the |proof|
% environment so that, when an optional title is given, the cross-reference it
% contains is turned into a |\proofof| (the title is typeset once into a
% discarded box with |\ref| and friends redefined). Without an optional title
% the original environment is used unchanged. |\protect| is forced to its
% typesetting meaning, so that a |\protect\ref| copied from a moving argument
% still reaches our redefinition; it must not be emptied instead, because the
% active characters of \textsf{babel} (the |~| of a title such as
% |[Proof of Theorem~\ref{thm:x}]|) expand through |\active@prefix|, which
% falls back on |\protect|\meta{char} when |\protect| holds neither of its
% two expected meanings -- with an empty |\protect| that reproduces the bare
% active character, endlessly.
%    \begin{macrocode}
\NewDocumentCommand\pgph@zcproofof{sO{}m}{\proofof{#3}}
\newcommand\pgph@scanproofarg[1]{%
  \begingroup
    \let\protect\@typeset@protect
    \def\ref##1{\proofof{##1}}\def\cref##1{\proofof{##1}}%
    \def\Cref##1{\proofof{##1}}\def\autoref##1{\proofof{##1}}%
    \let\zcref\pgph@zcproofof
    \setbox\z@\hbox{#1}%
  \endgroup
}
%    \end{macrocode}
% To capture the optional title of a detached |amsthm| proof
% (an optional |[of Theorem~\ref{thm:x}]|) we redefine the outer |\proof| so
% that, when a title is given, it is scanned and then handed to the
% \emph{saved} outer macro, which parses it as it always did. Delegating to the
% saved outer macro rather than to the inner one means we need to know neither
% the name of the inner macro nor the environment's default title, so the same
% wrapper serves any proof environment, whatever its default. No recursion
% occurs: the saved meaning is captured before the redefinition.
%
% We must wrap \emph{only} a genuine optional-argument environment: a package
% that captures the proof body verbatim (\apxproof, \textsf{myproofs}) redefines
% |proof| with no optional argument, and wrapping it would break that capture.
% We recognise the optional-argument form by the |\@protected@testopt| in its
% definition. When the proof environment is instead a theorem-like environment
% (Springer |\spnewtheorem*{proof}|, or a |\newtheorem| of one's own) its title
% is already handled by |\pgph@hookproofhead|, so there we just install the
% begin-hook.
%    \begin{macrocode}
\newcommand\pgph@wrapproofone[1]{%
  \@ifundefined{#1}{}{%
    \edef\pgph@proofmeaning{\expandafter\meaning\csname #1\endcsname}%
    \edef\pgph@testoptstr{\detokenize{\@protected@testopt}}%
    \IfSubStr\pgph@proofmeaning\pgph@testoptstr
      {\pgph@redefproof{#1}}%
      {\AtBeginEnvironment{#1}{\pgph@beginproof}}%
  }%
}
\newcommand\pgph@redefproof[1]{%
  \expandafter\let\csname pgph@savedenv@#1\expandafter\endcsname
    \csname #1\endcsname
  \expandafter\def\csname #1\endcsname{\pgph@beginproof
    \@ifnextchar[{\pgph@proofopt{#1}}%
                 {\csname pgph@savedenv@#1\endcsname}}%
}
\def\pgph@proofopt#1[#2]{%
  \pgph@scanproofarg{#2}\csname pgph@savedenv@#1\endcsname[{#2}]}
%    \end{macrocode}
% \end{macro}
% \end{macro}
%
% \DescribeMacro{\proofgraphproofenv}
% \begin{macro}{\proofgraphproofenv}
% \begin{macro}{\pgph@proofenvs}
% Which environments are proofs is a list, |proof| to begin with;
% |\proofgraphproofenv{|\meta{names}|}| adds to it, for a document whose proofs
% live in an environment of another name (\textsf{pf2}'s structured proofs under
% a name that avoids the clash with \textsf{amsthm}, a |\newtheorem{claimproof}|
% for the proofs of claims\dots). A proof is not a result, so the names are
% excluded from result tracking as well; |\proofgraphtrack| still overrides that
% for an environment one wants drawn as both.
%    \begin{macrocode}
\newcommand\pgph@proofenvs{proof}
\newcommand\proofgraphproofenv[1]{%
  \g@addto@macro\pgph@proofenvs{,#1}%
  \proofgraphuntrack{#1}%
}
%    \end{macrocode}
% \end{macro}
% \end{macro}
% \apxproof\ captures the body of a |proof| verbatim (to defer it to the
% appendix) and replays it there. Patching the |proof| environment then
% corrupts that capture, and the references are only executed at replay time
% anyway. But every proof \apxproof\ actually typesets runs through its saved
% copy of the original environment, |\axp@oldproof| -- a plain proof left in
% the body reaches it through the dispatcher, a proof replayed from the
% appendix file through the replay macros -- while the verbatim capture never
% invokes it. So under an \apxproof\ that fires the hook below we wrap that
% saved copy in place of |proof| itself, and likewise the |axp@old|\meta{name}
% copy of any other environment it defers (|claimproof|). Under an \apxproof\
% too old to fire the hook, wrapping would attribute each deferred proof to
% the result last typeset in the appendix instead of its theorem, so there we
% leave |proof| alone, as before. An environment added with
% |\proofgraphproofenv| that \apxproof\ does not defer is wrapped as usual.
% This happens at |\begin{document}|, once we know what is loaded and what
% the preamble has added to the list.
%    \begin{macrocode}
\newif\ifpgph@apxproof
\newif\ifpgph@apxhook
\newcommand\pgph@wrapproofs{%
  \@for\pgph@one:=\pgph@proofenvs\do{%
    \edef\pgph@pe{\expandafter\pgph@trimsp\expandafter{\pgph@one}}%
    \ifpgph@apxhook
      \ifcsname axp@old\pgph@pe\endcsname
        \edef\pgph@pe{axp@old\pgph@pe}%
      \fi
    \fi
    \ifboolexpr{ bool {pgph@apxproof} and test {\ifdefstring{\pgph@pe}{proof}} }%
      {}{\expandafter\pgph@wrapproofone\expandafter{\pgph@pe}}%
  }%
}
\AtBeginDocument{%
  \@ifpackageloaded{apxproof}{\pgph@apxprooftrue\pgph@apxinit}{}%
  \pgph@wrapproofs
}
%    \end{macrocode}
%
% \begin{macro}{\pgph@apxinit}
% \begin{macro}{\pgph@apxproofcapture}
% Under \apxproof, deferred proofs are typeset in the appendix, where their
% references finally execute. \apxproof\ fires |\apxproofhook| inside each
% such proof, passing the |axp@r|\meta{n} label of the proved theorem, which
% (because \apxproof\ also attaches that label to the theorem in the main
% text) is already in our map. An \apxproof\ that fires the hook also
% provides an empty default for it at load time, so finding it defined
% before we overwrite it with ours is how we know the hook comes
% (|\ifpgph@apxhook|, set here, read by |\pgph@wrapproofs| just after).
% When the hook fires, the wrapped |\axp@oldproof| has normally opened a
% proof instance already, and the label pins that instance, exactly as a
% |\proofof| in the proof would. The pin is expanded on the spot: what the
% hook passes is a macro holding the label, which the next deferred proof
% overwrites, so storing the token would resolve every pin to the last
% label of the run. When no instance is open (a proof
% environment the wrapping could not reach), we fall back to resolving the
% label to the source node and installing the reference-capturing commands
% directly, as before.
%    \begin{macrocode}
\newcommand\pgph@apxinit{%
  \ifdefined\apxproofhook\pgph@apxhooktrue\fi
  \def\apxproofhook##1{\pgph@apxproofcapture{##1}}}
\newcommand\pgph@apxproofcapture[1]{%
  \ifcsname pgph@powner@\pgph@curproof\endcsname
    \expandafter\xdef\csname pgph@pofl@\pgph@curproof\endcsname{#1}%
  \else
    \edef\pgph@curproof{\pgph@resolve{#1}}%
    \ifx\pgph@curproof\@empty\else
      \pgph@setpres\pgph@curproof
      \pgph@installcapture
    \fi
  \fi
}
%    \end{macrocode}
% \end{macro}
% \end{macro}
%
% \begin{macro}{\pgph@hookcite}
% Optional citation capture, applied to |\cite| and, when they are defined
% (\textsf{natbib}, \textsf{biblatex}), to |\citet| and |\citep| as well, all
% sharing one capture chain parameterized by the command name: external
% references become |cite:|-prefixed target
% labels, drawn later as distinct nodes. The node identity is the pair (key,
% note): a |\cite[|\meta{note}|]{|\meta{key}|}| inside a proof is distinguished
% from a |\cite| of the same key with a \emph{different} note, so that citing two
% different results of the same work (say two theorems of one book) yields two
% nodes rather than one mislabelled with whichever note was seen first. Each pair
% is assigned a numeric \emph{slot} (|\pgph@citeslot|); the edge target is
% |cite:|\meta{slot}, and the slot remembers the key (for the tag and the
% bibliography hyperlink) and the note (for the label). The note is reduced to a
% plain string at capture time (fragile wrappers neutralized, |~|~normalised to a
% space), which also serves as the signature so |[Thm 1]| and |[Thm~1]| coincide.
% We peek for the optional argument with |\@ifnextchar| rather than declaring
% one, so that a note-less |\cite{|\meta{key}|}| is passed on \emph{as written}:
% re-emitting it as |\cite[]{|\meta{key}|}| would make the typeset citation grow a
% spurious empty note (``[1, ]''). A second optional argument is peeked for in
% the same way, to support the two-argument form
% |\cite[|\meta{pre-note}|][|\meta{note}|]{|\meta{key}|}| of \textsf{natbib} and
% \textsf{biblatex}; there the \emph{second} argument is the note recorded for the
% node (matching the one-argument semantics, where the lone argument is the
% post-note), while the pre-note (``see'', ``e.g.'') qualifies the act of citing
% rather than identifying the cited result, and is only replayed to the saved
% |\cite|. Without this peek, the mandatory-argument scan would grab the bare
% token |[| as the key and typeset the rest of the citation as text, giving
% confusing errors (``Missing \$ inserted'') deep in the proof body. A star
% (|\citet*|, |\citep*|: full author list) is likewise peeked for and replayed;
% it affects rendering only, not the recorded (key, note) pair. As in
% |\pgph@hooknames|, the notes are replayed braced, |[{|\meta{note}|}]|, so that
% a note written |[{p.~37, [sic]}]| to hide a |]| reaches the saved command as
% written rather than cut short at that |]|. The replay head
% (saved command plus optional star) is frozen in |\pgph@replay| by the
% star-peek before the arguments are scanned. As in
% |\pgph@install|, installation is skipped (per command)
% when the command is already the capture macro; otherwise a proof nested inside
% another proof would save the capture as the saved command, and the first
% citation in the inner proof would recurse forever.
%    \begin{macrocode}
\newcommand\pgph@gxdefcs[2]{%
  \global\expandafter\edef\csname#1\endcsname{#2}}
\newcommand\pgph@cap@citegen[1]{%
  \def\pgph@citecmd{#1}%
  \@ifstar{\pgph@cap@cite@x*}{\pgph@cap@cite@x{}}}
\def\pgph@cap@cite@x#1{%
  \edef\pgph@replay{%
    \expandafter\noexpand\csname pgph@saved@\pgph@citecmd\endcsname#1}%
  \@ifnextchar[\pgph@cap@cite@opt\pgph@cap@cite@noopt}
\def\pgph@cap@cite@opt[#1]{%
  \@ifnextchar[{\pgph@cap@cite@optii{#1}}{\pgph@cap@cite@opti{#1}}}
\def\pgph@cap@cite@opti#1#2{%
  \pgph@cap@citerec{#1}{#2}\pgph@replay[{#1}]{#2}}
\def\pgph@cap@cite@optii#1[#2]#3{%
  \pgph@cap@citerec{#2}{#3}\pgph@replay[{#1}][{#2}]{#3}}
\def\pgph@cap@cite@noopt#1{\pgph@cap@citerec{}{#1}\pgph@replay{#1}}
\newcommand\pgph@cap@citerec[2]{%
  \ifpgph@nocapture\else
  \ifx\pgph@curproof\@empty\else
    \@for\pgph@one:=#2\do{%
      \edef\pgph@k{\expandafter\pgph@trimsp\expandafter{\pgph@one}}%
      \pgph@citeslot{\pgph@k}{#1}%
      \pgph@addedge{\pgph@curproof}{cite:\pgph@curslot}}%
  \fi\fi
}
\newcommand\pgph@citeslot[2]{%
  \begingroup
    \pgph@citesanitize
    \xdef\pgph@gnotestr{#2}%
  \endgroup
  \edef\pgph@sig{#1@@\detokenize\expandafter{\pgph@gnotestr}}%
  \ifcsname pgph@cslot@\pgph@sig\endcsname
    \edef\pgph@curslot{\csname pgph@cslot@\pgph@sig\endcsname}%
  \else
    \global\advance\pgph@citeslotcnt\@ne
    \edef\pgph@curslot{\the\pgph@citeslotcnt}%
    \pgph@gxdefcs{pgph@cslot@\pgph@sig}{\pgph@curslot}%
    \pgph@gxdefcs{pgph@cskey@\pgph@curslot}{#1}%
    \ifx\pgph@gnotestr\@empty\else
      \pgph@gxdefcs{pgph@citenote@\pgph@curslot}{\pgph@gnotestr}%
    \fi
  \fi
}
%    \end{macrocode}
% \end{macro}
%
% \begin{macro}{\pgph@citecmds}
% \begin{macro}{\proofgraphcitecommand}
% \begin{macro}{\pgph@installcite}
% Every citation command with the signature |\cs*[|\meta{pre}|][|\meta{post}|]|%
% \marg{keys} is captured by the same |\pgph@cap@citegen|, so which commands are
% captured is a list rather than code. The default list holds the citation
% commands of \LaTeX, \textsf{natbib} and \textsf{biblatex} that cite a work;
% |\proofgraphcitecommand| adds more. Naming a command that the document never
% defines costs nothing, as |\pgph@installciteone| installs only over an existing
% command, so the list may name every spelling a citation package might provide.
%
% The commands printing only a \emph{fragment} of a citation (|\citeauthor|,
% |\citeyear|\dots) are in the list as well. They commonly accompany a real
% citation of the same work, but a repetition is harmless: a (key, note) pair
% already seen reuses its slot, and |\pgph@writeedge| draws an edge once, so the
% second mention of a work adds nothing to the graph. Capturing them means that
% a proof invoking a work by author or by year alone, which does happen, is not
% silently left without its edge.
%
% The \textsf{biblatex} \emph{multicite} commands (|\cites|, |\parencites|,
% |\textcites|) cannot be captured at all: their
% |(|\meta{pre}|)(|\meta{post}|)| syntax and repeated key groups are a different
% grammar, which |\pgph@cap@citegen| does not parse.
%    \begin{macrocode}
\newcommand\pgph@citecmds{%
  cite,Cite,citet,Citet,citep,Citep,%
  citealt,Citealt,citealp,Citealp,%
  parencite,Parencite,textcite,Textcite,autocite,Autocite,%
  footcite,Footcite,footcitetext,smartcite,Smartcite,supercite,%
  fullcite,footfullcite,%
  citeauthor,Citeauthor,citefullauthor,citetitle,Citetitle,%
  citeyear,Citeyear,citeyearpar,citedate,Citedate,citeurl,citenum}
\newcommand\proofgraphcitecommand[1]{\g@addto@macro\pgph@citecmds{,#1}}
\newcommand\pgph@installcite{%
  \@for\pgph@conename:=\pgph@citecmds\do{%
    \edef\pgph@cone{\expandafter\pgph@trimsp\expandafter{\pgph@conename}}%
    \expandafter\pgph@installciteone\expandafter{\pgph@cone}}%
}
%    \end{macrocode}
% The capture macro of a command is built on first use rather than declared: it
% is the same |\pgph@cap@citegen| call for all of them, and building it only for
% commands the document actually defines keeps the unused names of the list from
% costing anything. It must be |\gdef|, the installation happening inside the
% group of a proof.
%    \begin{macrocode}
\newcommand\pgph@installciteone[1]{%
  \ifcsname #1\endcsname
    \ifcsname pgph@cap@#1\endcsname\else
      \expandafter\gdef\csname pgph@cap@#1\endcsname{\pgph@cap@citegen{#1}}%
    \fi
    \expandafter\ifx\csname #1\expandafter\endcsname
                    \csname pgph@cap@#1\endcsname
    \else
      \expandafter\let\csname pgph@saved@#1\expandafter\endcsname
                      \csname #1\endcsname
      \expandafter\let\csname #1\expandafter\endcsname
                      \csname pgph@cap@#1\endcsname
    \fi
  \fi
}
%    \end{macrocode}
% \end{macro}
% \end{macro}
% \end{macro}
%
% The \textsf{biblatex} recorder is registered at the end of the preamble, once
% \textsf{biblatex} has certainly been loaded (it is a preamble-only package)
% but while |\AtEveryCitekey| still accepts additions. It runs both at every
% citation and at every bibliography entry, so that a work only |\nocite|d, or
% reached by |\usescite| alone, is still seen. It is not conditioned on the
% |cite| option, which governs the automatic capture of |\cite| only:
% |\usescite| and |\proofgraphciteedge| draw citation nodes without it, and
% those nodes want their tag and their anchor just the same.
%    \begin{macrocode}
\newif\ifpgph@blx
\AtEndPreamble{%
  \@ifpackageloaded{biblatex}{%
    \pgph@blxtrue
    \AtEveryCitekey{\pgph@blxrecord}%
    \AtEveryBibitem{\pgph@blxrecord}%
  }{}%
}
%    \end{macrocode}
%
% \subsection{User commands}
%    \begin{macrocode}
\newcommand\uses[1]{\pgph@recordref{#1}}
%    \end{macrocode}
% |\proofof| names the result a proof establishes. Inside a proof instance it
% records the label against the instance, so that the pin covers the whole
% proof, wherever in it the command stands, and is resolved (with every other
% label) when the graph is written. Outside one -- in an \apxproof\ deferred
% proof, whose owning result the |\apxproofhook| already gives, or in the body
% of a result under |statements| -- there is no instance to pin, and the source
% is redirected on the spot instead.
%    \begin{macrocode}
\newcommand\proofof[1]{%
  \ifcsname pgph@powner@\pgph@curproof\endcsname
    \expandafter\gdef\csname pgph@pofl@\pgph@curproof\endcsname{#1}%
  \else
    \ifcsname pgph@map@#1\endcsname
      \xdef\pgph@curproof{\csname pgph@map@#1\endcsname}%
    \fi
  \fi
}
\newcommand\proofgraphedge[2]{%
  \g@addto@macro\pgph@edges{\pgph@labeledge{#1}{#2}}%
}
\newcommand\pgph@labeledge[2]{%
  \edef\pgph@ef{\pgph@resolve{#1}}%
  \ifx\pgph@ef\@empty\else
    \@for\pgph@one:=#2\do{%
      \edef\pgph@k{\expandafter\pgph@trimsp\expandafter{\pgph@one}}%
      \pgph@plainedge{\pgph@ef}{\pgph@k}}%
  \fi
}
%    \end{macrocode}
% \DescribeMacro{\usescite}
% \DescribeMacro{\proofgraphciteedge}
% \begin{macro}{\usescite}
% \begin{macro}{\proofgraphciteedge}
% \begin{macro}{\pgph@citelabeledge}
% The two manual counterparts of the automatic |\cite| capture, for a dependency
% on external work that no citation in a proof expresses: |\usescite|, the
% |\uses| of citation nodes, records one from inside a proof, and
% |\proofgraphciteedge|, the |\proofgraphedge| of citation nodes, states one
% from anywhere. Both create the citation node if the work has no node yet, take
% a comma-separated list of \BibTeX\ keys and accept as an optional argument the
% note that would have been the optional argument of |\cite| (|[Thm~3]|), which
% labels the node and, as in the automatic capture, distinguishes one result of
% a work from another.
%
% Being explicit, they do not depend on the |cite| option, which only says
% whether a plain |\cite| is captured on its own. |\usescite| shares the
% recording path of the captured citation commands, so a work reached both ways
% still yields a single node.
%    \begin{macrocode}
\NewDocumentCommand\usescite{O{}m}{\pgph@cap@citerec{#1}{#2}}
\NewDocumentCommand\proofgraphciteedge{O{}mm}{%
  \g@addto@macro\pgph@edges{\pgph@citelabeledge{#1}{#2}{#3}}%
}
%    \end{macrocode}
% Resolution is deferred to emission, as for |\proofgraphedge|, so the result
% may be labelled anywhere in the document; the slot is therefore allocated
% then, which is harmless as slot numbers only have to be distinct. The edge is
% assembled by |\edef| rather than passed as macros, because |\pgph@citeedge|
% takes the |cite:|\meta{slot} target apart with \textsf{xstring}, which reads
% its argument literally.
%    \begin{macrocode}
\newcommand\pgph@citelabeledge[3]{%
  \edef\pgph@ef{\pgph@resolve{#2}}%
  \ifx\pgph@ef\@empty\else
    \@for\pgph@one:=#3\do{%
      \edef\pgph@k{\expandafter\pgph@trimsp\expandafter{\pgph@one}}%
      \pgph@citeslot{\pgph@k}{#1}%
      \edef\pgph@tmpce{%
        \noexpand\pgph@citeedge{\pgph@ef}{cite:\pgph@curslot}}%
      \pgph@tmpce}%
  \fi
}
%    \end{macrocode}
% \end{macro}
% \end{macro}
% \end{macro}
% \begin{macro}{\proofgraphignore}
% \begin{macro}{\proofgraphignorecite}
% \begin{macro}{\proofgraphexclude}
% The editing rules. Each undoes one of the recording commands and takes the
% arguments of the command it undoes: |\proofgraphignore| those of
% |\proofgraphedge|, the result first and what its proof uses second, and
% |\proofgraphignorecite| those of |\proofgraphciteedge|, the result and then
% the keys, with the note of the citation as an optional argument. So a rule is
% written by copying the dependency it is to remove, not by reading the arrow
% off the picture, which with the default |direction=usedby| runs the other way.
% The |direction| option is not consulted here at all: the filters work on the
% relation recorded, and |direction| settles only how it is drawn.
%    \begin{macrocode}
\newcommand\proofgraphignore[2]{%
  \g@addto@macro\pgph@ignores{\pgph@ignorerule{#1}{#2}}%
}
\NewDocumentCommand\proofgraphignorecite{O{}mm}{%
  \g@addto@macro\pgph@ignores{\pgph@citeignorerule{#1}{#2}{#3}}%
}
\newcommand\proofgraphexclude[1]{%
  \g@addto@macro\pgph@excludes{\pgph@excluderule{#1}}%
}
%    \end{macrocode}
% \end{macro}
% \end{macro}
% \end{macro}
% The attributes go to \Graphviz\ untouched, and a colour is written there as
% |"#ff0000"|. A |#| read with its usual category is a macro parameter
% character, which no macro may simply hold: the style is kept in a macro, so
% such an attribute list stopped the run with ``Illegal parameter number'',
% first as the style was declared and again as it was written out. Doubling the
% |#| does not help either, the second error remaining and |\write| doubling it
% back in the |.dot|. The attributes are therefore read with |#| made an
% ordinary character, which passes through the macro, the |\edef| assembling
% the attribute list of a node and the write to the |.dot| as itself. Reading
% them takes two steps: the category has to be changed before the argument is
% read, so the command takes no argument of its own and hands over to a
% delimiter-free helper.
%    \begin{macrocode}
\newcommand\proofgraphstyle{\begingroup\@makeother\#\pgph@setstyle}
\newcommand\pgph@setstyle[2]{\endgroup
  \global\@namedef{pgph@style@#1}{#2}}
%    \end{macrocode}
% External citation nodes (option |cite|) share one style, held in
% |\pgph@citestyle| and overridable with |\proofgraphstylecite|; it defaults to
% |shape=note|.
%    \begin{macrocode}
\newcommand\pgph@citestyle{shape=note}
\newcommand\proofgraphstylecite{\begingroup\@makeother\#\pgph@setcitestyle}
\newcommand\pgph@setcitestyle[1]{\endgroup
  \renewcommand\pgph@citestyle{#1}}
%    \end{macrocode}
%
% \subsection{Resolution helpers}
% These run at end of document, after the |.aux| has populated the map.
%    \begin{macrocode}
\newcommand\pgph@resolve[1]{%
  \ifcsname pgph@map@#1\endcsname\csname pgph@map@#1\endcsname\fi
}
% The key under which a filter is consulted is that of the relation recorded,
% whose source is the \emph{result}; the two names below say which end is which,
% so that reading a rule is not a matter of remembering an order.
%    \begin{macrocode}
\newcommand\pgph@ignorerule[2]{%
  \edef\pgph@ires{\pgph@resolve{#1}}\edef\pgph@idep{\pgph@resolve{#2}}%
  \ifx\pgph@ires\@empty\else\ifx\pgph@idep\@empty\else
    \@namedef{pgph@ig@\pgph@ires @\pgph@idep}{}%
  \fi\fi
}
%    \end{macrocode}
% A citation rule names the work depended on by key and note, the pair a
% citation node stands for, so it has to go through the slot allocator to learn
% which node that is. Asking for the slot of a work no proof cites allocates
% one, which costs nothing: slots become nodes only when an edge reaches them.
%    \begin{macrocode}
\newcommand\pgph@citeignorerule[3]{%
  \edef\pgph@ires{\pgph@resolve{#2}}%
  \ifx\pgph@ires\@empty\else
    \@for\pgph@one:=#3\do{%
      \edef\pgph@k{\expandafter\pgph@trimsp\expandafter{\pgph@one}}%
      \pgph@citeslot{\pgph@k}{#1}%
      \@namedef{pgph@igc@\pgph@ires @\pgph@curslot}{}}%
  \fi
}
\newcommand\pgph@excluderule[1]{%
  \edef\pgph@x{\pgph@resolve{#1}}%
  \ifx\pgph@x\@empty\else\@namedef{pgph@ex@\pgph@x}{}\fi
}
%    \end{macrocode}
% \begin{macro}{\pgph@checkrules}
% An editing rule that matches nothing does nothing, and says nothing: the
% commonest way to get one is to write the two labels of |\proofgraphignore| in
% the order of the arrow drawn rather than of the dependency recorded, which
% leaves the graph exactly as it was, with no sign that a rule was even read. So
% the rules are replayed once the edges have been written, against the relations
% the document turned out to record, and each rule that met none is reported.
% When it is the swapped rule that would have matched, the report says so and
% gives the line to write, that being the mistake it almost always is.
%
% The replay redefines the three rule macros and runs the same token lists
% again, so a rule cannot be checked in a way that differs from how it was
% applied, and adding a kind of rule cannot leave the check behind.
%    \begin{macrocode}
\newcommand\pgph@checkrules{%
  \let\pgph@ignorerule\pgph@ignorecheck
  \let\pgph@citeignorerule\pgph@citeignorecheck
  \let\pgph@excluderule\pgph@excludecheck
  \pgph@ignores
  \pgph@excludes
}
\newcommand\pgph@nolabelwarn[2]{%
  \PackageWarning{proofgraph}{%
    No such result in the graph:\MessageBreak
    \string#1\space names `#2', which is not\MessageBreak
    the label of any result drawn}%
}
\newcommand\pgph@ignorecheck[2]{%
  \edef\pgph@ires{\pgph@resolve{#1}}\edef\pgph@idep{\pgph@resolve{#2}}%
  \ifx\pgph@ires\@empty
    \pgph@nolabelwarn\proofgraphignore{#1}%
  \else\ifx\pgph@idep\@empty
    \pgph@nolabelwarn\proofgraphignore{#2}%
  \else\ifcsname pgph@rel@\pgph@ires @\pgph@idep\endcsname\else
    \ifcsname pgph@rel@\pgph@idep @\pgph@ires\endcsname
      \def\pgph@hint{%
        it is `#2' that depends on `#1'.\MessageBreak
        The two labels go the way the dependency is\MessageBreak
        recorded, the result first, and not the way\MessageBreak
        the arrow is drawn. Write instead\MessageBreak
        \string\proofgraphignore{#2}{#1}}%
    \else
      \def\pgph@hint{%
        nothing records that `#1' depends on `#2'}%
    \fi
    \PackageWarning{proofgraph}{%
      No edge suppressed by\MessageBreak
      \string\proofgraphignore{#1}{#2}:\MessageBreak
      \pgph@hint}%
  \fi\fi\fi
}
\newcommand\pgph@citeignorecheck[3]{%
  \edef\pgph@ires{\pgph@resolve{#2}}%
  \ifx\pgph@ires\@empty
    \pgph@nolabelwarn\proofgraphignorecite{#2}%
  \else
    \edef\pgph@tmpnote{\detokenize{#1}}%
    \ifx\pgph@tmpnote\@empty
      \def\pgph@hint{}%
    \else
      \def\pgph@hint{\space with the note given}%
    \fi
    \@for\pgph@one:=#3\do{%
      \edef\pgph@k{\expandafter\pgph@trimsp\expandafter{\pgph@one}}%
      \pgph@citeslot{\pgph@k}{#1}%
      \ifcsname pgph@relc@\pgph@ires @\pgph@curslot\endcsname\else
        \PackageWarning{proofgraph}{%
          No edge suppressed by\MessageBreak
          \string\proofgraphignorecite{#2}{\pgph@k}:\MessageBreak
          nothing records that `#2' cites that work\pgph@hint}%
      \fi}%
  \fi
}
\newcommand\pgph@excludecheck[1]{%
  \edef\pgph@x{\pgph@resolve{#1}}%
  \ifx\pgph@x\@empty\pgph@nolabelwarn\proofgraphexclude{#1}\fi
}
%    \end{macrocode}
% \end{macro}
%
% \subsection{Emission}
% \begin{macro}{\pgph@dotesc}
% Escape the two characters that are special in a \Graphviz\ double-quoted string,
% so a theorem name, number or citation note containing them does not corrupt the
% |.dot|: a backslash is doubled and a double quote is backslash-escaped (in that
% order). |\pgph@bschar| and |\pgph@dqchar| hold the catcode-12 backslash and
% double quote to search for; \textsf{xstring}'s expansion mode is saved and
% restored around the two substitutions.
%    \begin{macrocode}
\begingroup
  \catcode`\|=0 \catcode`\"=12 \catcode`\\=12
  |gdef|pgph@bschar{\}%
  |gdef|pgph@dqchar{"}%
|endgroup
\newcommand\pgph@dotesc[1]{%
  \saveexpandmode\expandarg
  \StrSubstitute{#1}{\pgph@bschar}{\pgph@bschar\pgph@bschar}[#1]%
  \StrSubstitute{#1}{\pgph@dqchar}{\pgph@bschar\pgph@dqchar}[#1]%
  \restoreexpandmode}
%    \end{macrocode}
% \end{macro}
% \begin{macro}{\pgph@node}
% Emit one node line, honouring exclusions and per-environment styling.
% External (|cite:|) targets are emitted separately as a pass over edges.
% An \emph{unnumbered} result (no number captured) is omitted: a results graph
% is about numbered statements, and unnumbered theorem-like environments are
% typically expository or repeated copies (e.g., a class-predefined
% |\spnewtheorem*{claim}|). |\proofgraphtrack{|\meta{env}|}| overrides this, so a
% deliberately unnumbered environment is still drawn.
%    \begin{macrocode}
\newcommand\pgph@node[3]{%
  \ifcsname pgph@ex@#1\endcsname\else
    \edef\pgph@nm{\ifcsname pgph@num@#1\endcsname%
      \csname pgph@num@#1\endcsname\fi}%
    \ifx\pgph@nm\@empty
      \ifcsname pgph@force@#2\endcsname
        \pgph@emitnode{#1}{#2}{#3}%
      \else
        \@namedef{pgph@ex@#1}{}%
      \fi
    \else
      \pgph@emitnode{#1}{#2}{#3}%
    \fi
  \fi
}
\newcommand\pgph@emitnode[3]{%
  \edef\pgph@s{\ifcsname pgph@style@#2\endcsname
    , \csname pgph@style@#2\endcsname\fi}%
  \edef\pgph@dn{\ifcsname pgph@name@#1\endcsname
    \csname pgph@name@#1\endcsname\else#3\fi}%
  \edef\pgph@lbl{\pgph@dn\space\pgph@nm}%
  \pgph@dotesc\pgph@lbl
  \immediate\write\pgph@out{%
    \space\space"n#1" [label="\pgph@lbl"\pgph@s];}%
  \ifcsname pgph@rmap@#1\endcsname
    \immediate\write\pgph@mapout{%
      \string\pgphnodemap{n#1}{\csname pgph@rmap@#1\endcsname}}%
  \fi
}
%    \end{macrocode}
% \end{macro}
%
% \begin{macro}{\pgph@edge}
% \begin{macro}{\pgph@setsrc}
% Emit one edge, resolving both ends and applying the self-loop, ignore and
% exclude filters. Targets beginning with |cite:| denote external nodes.
%
% The source is resolved by |\pgph@setsrc|. An edge captured in a proof has that
% proof instance as its source, and the result it belongs to is decided here:
% the label |\proofof| pinned the proof to, when there is one and it resolves,
% and otherwise the result the proof followed. An unresolvable pin falls back on
% that result rather than dropping the edge, so a mistyped or never-labelled
% |\proofof| loses the attribution but not the dependency. An edge recorded
% outside a proof (the body or the title of a result, an \apxproof\ deferred
% proof) already carries a result id, which is its own source.
%
% A sub-proof owns not a result but the proof it is nested in, so the owner is
% followed as a chain: at each step the pin of the instance reached wins if it
% resolves, and otherwise the walk goes on to its owner. The nearest pin
% therefore decides, an inner |\proofof| overriding an outer one for the
% references of the inner proof alone, and an outer one covering every nested
% reference no inner pin claims. The walk always terminates: an owner is either
% an instance opened strictly earlier or a result id, which owns nothing.
%    \begin{macrocode}
\newcommand\pgph@setsrc[1]{%
  \edef\pgph@esrc{#1}%
  \pgph@srcwalk
}
\newcommand\pgph@srcwalk{%
  \ifcsname pgph@powner@\pgph@esrc\endcsname
    \let\pgph@srcnext\pgph@srcstep
  \else
    \let\pgph@srcnext\relax
  \fi
  \pgph@srcnext
}
\newcommand\pgph@srcstep{%
  \let\pgph@epin\@empty
  \ifcsname pgph@pofl@\pgph@esrc\endcsname
    \edef\pgph@epin{\csname pgph@pofl@\pgph@esrc\endcsname}%
    \edef\pgph@epin{\pgph@resolve{\pgph@epin}}%
  \fi
  \ifx\pgph@epin\@empty
    \edef\pgph@esrc{\csname pgph@powner@\pgph@esrc\endcsname}%
    \pgph@srcwalk
  \else
    \let\pgph@esrc\pgph@epin
  \fi
}
\newcommand\pgph@edge[2]{%
  \pgph@setsrc{#1}%
  \ifx\pgph@esrc\@empty\else
    \expandafter\pgph@edgeone\expandafter{\pgph@esrc}{#2}%
  \fi
}
\newcommand\pgph@edgeone[2]{%
  \IfBeginWith{#2}{cite:}%
    {\pgph@citeedge{#1}{#2}}%
    {\IfBeginWith{#2}{node:}%
      {\StrBehind{#2}{node:}[\pgph@ntgt]%
       \pgph@plainedgeid{#1}{\pgph@ntgt}}%
      {\pgph@plainedge{#1}{#2}}}%
}
%    \end{macrocode}
% A target is normally a user label, resolved here; a |node:| one is a result id
% already, recorded by |\pgph@supportedge|, which knows the result it points at
% without its having to be labelled at all. Past that, the two are the same
% edge and take the same filters.
%    \begin{macrocode}
\newcommand\pgph@plainedge[2]{%
  \edef\pgph@ptgt{\pgph@resolve{#2}}%
  \ifx\pgph@ptgt\@empty\else
    \pgph@plainedgeid{#1}{\pgph@ptgt}%
  \fi
}
\newcommand\pgph@plainedgeid[2]{%
  \edef\pgph@t{#2}%
  \@namedef{pgph@rel@#1@\pgph@t}{}%
  \ifcsname pgph@ex@#1\endcsname\else
  \ifcsname pgph@ex@\pgph@t\endcsname\else
  \ifcsname pgph@ig@#1@\pgph@t\endcsname\else
    \ifboolexpr{ test {\ifnumequal{#1}{\pgph@t}}
                 and test {\ifdefstring{\pgph@selfloops}{remove}} }%
      {}{\pgph@emitedge{n#1}{n\pgph@t}}%
  \fi\fi\fi
}
%    \end{macrocode}
% \end{macro}
% \end{macro}
%
% \begin{macro}{\pgph@emitedge}
% \begin{macro}{\pgph@writeedge}
% The relation recorded is always ``the result whose proof we are in (the first
% argument) uses the target (the second)''. The |direction| option chooses the
% arrow orientation: |direction=usedby| (default) draws \emph{target to result},
% so an arrow \texttt{a->b} reads ``\texttt{a} is used by \texttt{b}'' (i.e.\
% \texttt{b} relies on \texttt{a}); |direction=uses| draws \emph{result to
% target} (``\texttt{a} uses \texttt{b}''). The filters above work on the
% recorded relation, so they are unaffected by the chosen orientation.
% |\pgph@writeedge| deduplicates.
%    \begin{macrocode}
\newcommand\pgph@emitedge[2]{%
  \ifdefstring{\pgph@direction}{uses}%
    {\pgph@writeedge{#1}{#2}}%
    {\pgph@writeedge{#2}{#1}}%
}
\newcommand\pgph@writeedge[2]{%
  \ifcsname pgph@drawn@#1@#2\endcsname\else
    \@namedef{pgph@drawn@#1@#2}{}%
    \immediate\write\pgph@out{\space\space"#1" -> "#2";}%
  \fi
}
%    \end{macrocode}
% \end{macro}
% \end{macro}
%
% \begin{macro}{\pgph@citeedge}
% \begin{macro}{\pgph@setcitelabel}
% A |cite:|\meta{slot} target: recover the slot's \BibTeX\ key, declare the
% external node once per slot, then draw the edge. The node \emph{label} is built
% by |\pgph@setcitelabel|: it is the key by default, or, with |citelabel=tag|, the
% printed citation tag (the ``[42]''/``[Knu74]'' text), recovered either from a
% tag recorded by |\pgph@blxrecord| under \textsf{biblatex} or from the
% |\b@|\meta{key} macro that |\bibcite| writes to the |.aux| (so two runs are
% needed, and it falls back to the key when no tag is known, e.g., on the first
% run or with citation packages neither path covers). The slot's
% |\cite| note, when present, is added inside the tag. The
% node-to-label map records the \emph{anchor} of the cited work, so all slots of
% one work hyperlink to the same bibliography entry.
%    \begin{macrocode}
\newcommand\pgph@citeedge[2]{%
  \StrBehind{#2}{cite:}[\pgph@slot]%
  \@namedef{pgph@relc@#1@\pgph@slot}{}%
  \ifcsname pgph@igc@#1@\pgph@slot\endcsname\else
  \ifcsname pgph@ex@#1\endcsname\else
    \edef\pgph@key{\csname pgph@cskey@\pgph@slot\endcsname}%
    \ifcsname pgph@extnode@\pgph@slot\endcsname\else
      \@namedef{pgph@extnode@\pgph@slot}{}%
      \expandafter\pgph@setcitelabel\expandafter{\pgph@slot}%
      \pgph@dotesc\pgph@clabel
      \immediate\write\pgph@out{%
        \space\space"c@\pgph@slot"
        [label="\pgph@clabel",\pgph@citestyle];}%
      \immediate\write\pgph@mapout{%
        \string\pgphnodemapcite{c@\pgph@slot}{\pgph@citeanchor{\pgph@key}}}%
    \fi
    \pgph@emitedge{n#1}{c@\pgph@slot}%
  \fi\fi
}
\edef\pgph@obrace{\expandafter\@gobble\string\{}
\newif\ifpgph@natnum
\newcommand\pgph@nil{}
%    \end{macrocode}
% |\pgph@citesanitize| neutralizes the wrappers and spacing cruft that citation
% labels carry, so that an |\edef| of |\b@|\meta{key} yields plain text:
% \textsf{hyperref}'s |\hyper@@link| (keep its printed last argument) and the
% boxing/penalty/|\ignorespaces| noise that \textsf{natbib} bakes into author
% lists.
%    \begin{macrocode}
\newcommand\pgph@citesanitize{%
  \let\protect\@empty
  \def\hyper@@link[##1]##2##3##4{##4}%
  \let\unskip\@empty \let\ignorespaces\@empty \let\unhbox\@empty
  \let\penalty\@gobble \let\@M\@empty \let\voidb@x\@empty
  \let\nobreakspace\space \let~\space
  \def\hbox##1{##1}\def\mbox##1{##1}%
}
%    \end{macrocode}
% A \textsf{natbib} |\b@|\meta{key} is |{|\meta{num}|}{|\meta{year}|}{|%
% \meta{short}|}{|\meta{long}|}|. In numbers mode the tag is the number; in
% author--year mode it is the short author list and the year.
%    \begin{macrocode}
\def\pgph@natpick#1#2#3#4\pgph@nil{%
  \ifpgph@natnum\def\pgph@x{#1}\else\def\pgph@x{#3 #2}\fi}
\newcommand\pgph@natdo{\expandafter\pgph@natpick\pgph@x{}{}{}\pgph@nil}
%    \end{macrocode}
% \begin{macro}{\pgph@natnumset}
% Which of the two \textsf{natbib} modes is in force is read off
% |\ifNAT@numbers|, by comparing it with |\iftrue|. That comparison
% \emph{must} live in a macro of its own rather than inline in
% |\pgph@setcitelabel|, because the code around it can be skipped by a false
% conditional: skipping does not expand anything, it merely counts conditional
% tokens, and the |\iftrue| here -- an operand of |\ifx|, not a conditional --
% would be counted as one more |\if| to be closed. The skip would then eat one
% |\fi| too many and run past the end of |\pgph@setcitelabel| into its caller,
% silently swallowing the very writes that emit the citation node. Inside a
% macro the token is never seen by the skipping scanner.
%    \begin{macrocode}
\newcommand\pgph@natnumset{%
  \pgph@natnumfalse
  \@ifundefined{ifNAT@numbers}{}{%
    \expandafter\ifx\csname ifNAT@numbers\endcsname\iftrue
      \pgph@natnumtrue\fi}%
}
%    \end{macrocode}
% \end{macro}
%
% \begin{macro}{\pgph@blxrecord}
% \textsf{biblatex} writes no |\b@|\meta{key}: its tags are assembled while the
% citation is typeset, from data the |.bbl| holds. We therefore record them as
% they go by, from |\AtEveryCitekey| and |\AtEveryBibitem| (installed below),
% where the entry is the current one and its fields are readable. |labelnumber|
% carries the tag of the
% numeric styles and |labelalpha| that of the alphabetic ones; the author--year
% styles have neither, as their tag is a formatted name list rather than a
% field, and are left to the |key| fallback. A field that survives sanitizing
% with a backslash still in it is not plain text, and is likewise declined.
%
% The same visit records the entry's \emph{anchor}, the name of the
% \textsf{hyperref} destination its bibliography entry carries, which
% \textsf{biblatex} does not spell |cite.|\meta{key} as the \LaTeX\ kernel and
% \textsf{natbib} do but |cite.|\meta{refsection}|@|\meta{key}, the reference
% section being part of the name because one key may have an entry in each.
% Naming the destination wrongly is not a visible failure but a misleading one:
% \textsf{hyperref} sends an unknown destination to the end of the document, so
% the node would link somewhere rather than nowhere. The anchor is recorded
% whatever |citelabel| says, the hyperlink being drawn either way.
%    \begin{macrocode}
\newcommand\pgph@blxrecord{%
  \begingroup
    \pgph@citesanitize
    \edef\pgph@blxkey{\thefield{entrykey}}%
    \pgph@gxdefcs{pgph@blxanchor@\pgph@blxkey}%
      {\the\c@refsection @\pgph@blxkey}%
    \ifdefstring{\pgph@citelabel}{tag}{%
      \edef\pgph@x{\thefield{labelnumber}}%
      \ifx\pgph@x\@empty \edef\pgph@x{\thefield{labelalpha}}\fi
      \edef\pgph@dx{\detokenize\expandafter{\pgph@x}}%
      \ifx\pgph@x\@empty\else
        \IfSubStr\pgph@dx\@backslashchar{}{%
          \pgph@gxdefcs{pgph@blxtag@\pgph@blxkey}{\pgph@x}}%
      \fi
    }{}%
  \endgroup
}
%    \end{macrocode}
% The anchor of a key never met as a citation nor as a bibliography entry is not
% known; under \textsf{biblatex} we then assume the first reference section,
% which is the only one most documents have. A key absent from the bibliography
% has no destination under any name, so nothing is lost when the guess is wrong.
%    \begin{macrocode}
\newcommand\pgph@citeanchor[1]{%
  \ifcsname pgph@blxanchor@#1\endcsname
    \csname pgph@blxanchor@#1\endcsname
  \else
    \ifpgph@blx 0@\fi#1%
  \fi
}
%    \end{macrocode}
% \end{macro}
%
% A note is written the way a citation with a note is printed, inside the tag
% and after it: |[Knu68, pp. 23-27]|, not the tag after the note. Its dashes are
% the two ligatures \LaTeX\ reads as dashes, |--| and |---|, which mean nothing
% to \Graphviz; they are turned into the \textsc{html} character entities
% \Graphviz\ decodes in an ordinary label, so that a page range set as
% |[pp.~23--27]| carries the same dash in the graph as in the document. The
% entities are built with |&| of category~12, as the escapes of
% |\pgph@dotesc| are, so that nothing in the label depends on the category codes
% in force where the graph is written.
%    \begin{macrocode}
\begingroup
  \catcode`\&=12
  \gdef\pgph@endash{&ndash;}%
  \gdef\pgph@emdash{&mdash;}%
\endgroup
\newcommand\pgph@dashesc[1]{%
  \saveexpandmode\expandarg
  \StrSubstitute{#1}{---}{\pgph@emdash}[#1]%
  \StrSubstitute{#1}{--}{\pgph@endash}[#1]%
  \restoreexpandmode}
%    \end{macrocode}
%    \begin{macrocode}
\newif\ifpgph@ctag
\newcommand\pgph@setcitelabel[1]{%
  \edef\pgph@ckey{\csname pgph@cskey@#1\endcsname}%
  \edef\pgph@cbase{\pgph@ckey}%
  \pgph@ctagfalse
  \ifdefstring{\pgph@citelabel}{tag}{%
    \ifcsname pgph@blxtag@\pgph@ckey\endcsname
      \edef\pgph@cbase{\csname pgph@blxtag@\pgph@ckey\endcsname}%
      \pgph@ctagtrue
    \else
    \ifcsname b@\pgph@ckey\endcsname
      \global\let\pgph@cbaseg\relax
      \begingroup
        \pgph@citesanitize
        \edef\pgph@x{\csname b@\pgph@ckey\endcsname}%
        \edef\pgph@dx{\detokenize\expandafter{\pgph@x}}%
        \IfBeginWith\pgph@dx\pgph@obrace{%
          \pgph@natnumset
          \pgph@natdo
          \edef\pgph@dx{\detokenize\expandafter{\pgph@x}}%
        }{}%
        \ifx\pgph@x\@empty\else
          \IfSubStr\pgph@dx\@backslashchar{}{%
            \global\let\pgph@cbaseg\pgph@x}%
        \fi
      \endgroup
      \ifx\pgph@cbaseg\relax\else
        \let\pgph@cbase\pgph@cbaseg
        \pgph@ctagtrue
      \fi
    \fi\fi
  }{}%
  \edef\pgph@cnote{%
    \ifcsname pgph@citenote@#1\endcsname
      , \csname pgph@citenote@#1\endcsname\fi}%
  \ifx\pgph@cnote\@empty\else\pgph@dashesc\pgph@cnote\fi
  \edef\pgph@clabel{%
    \ifpgph@ctag[\fi\pgph@cbase\pgph@cnote\ifpgph@ctag]\fi}%
}
%    \end{macrocode}
% \end{macro}
% \end{macro}
%
% \begin{macro}{\pgph@render}
% \begin{macro}{\pgph@checkrender}
% \begin{macro}{\pgph@maybrender}
% \begin{macro}{\pgph@lastrender}
% |autorun| must not re-render a graph that has not changed. Rendering
% unconditionally is not merely wasteful: the image \Graphviz\ produces is not
% guaranteed to be reproducible (its \textsf{cairo}-based outputs, |pdf| among
% them, carry a creation date), so an unconditional run makes the embedded image
% differ at every compilation. A build tool such as \textsf{latexmk} then sees an
% input file that never settles and reruns until it gives up, reporting that too
% many passes were needed.
%
% We therefore render only when something relevant has changed. The signature of
% a rendering is the checksum of the |.dot| file just written, together with the
% engine and the format, which the |.dot| does not record; it is stored in the
% |.aux| by |\pgph@lastrender| and compared with the signature of the next run.
% \Graphviz\ is run when the signature differs, when no signature is known yet
% (first run), when either output file is missing, and, as a safe fallback, when
% the engine provides no file checksum.
%
% Both signatures are |\detokenize|d before being compared. The checksum comes
% out of |\pdf@filemdfivesum| as characters of category \emph{other}, while the
% same string read back from the |.aux| has its letters as category
% \emph{letter}, so an |\ifx| on the raw strings would never see two runs as
% equal.
%    \begin{macrocode}
\let\pgph@prevsig\relax
\newcommand\pgph@lastrender[1]{\xdef\pgph@prevsig{\detokenize{#1}}}
\newcommand\pgph@render{%
  \ShellEscape{\pgph@engine\space -T\pgph@format\space
    \pgph@file.dot -o \pgph@file-proofgraph.\pgph@format}%
  \ShellEscape{\pgph@engine\space -Tplain\space
    \pgph@file.dot -o \pgph@file-proofgraph.plain}%
  \pgph@checkrender
}
%    \end{macrocode}
% A shell escape that does not happen says nothing: |\ShellEscape| reports no
% failure, and the run continues as if \Graphviz\ had been called. The
% rendering is therefore checked rather than assumed, right after the command,
% where the file it was to produce is already visible to |\IfFileExists|. What
% the user would otherwise meet is the warning of |\proofgraph| on the
% \emph{next} run, saying only that the image is not there.
%
% |\ShellEscapeStatus| tells the three cases apart: no shell escape at all, an
% unrestricted one (so the run was attempted and it is \Graphviz\ that failed
% or is missing), and the restricted mode a \TeX\ installation has by default,
% which runs only the commands the installation lists and never \Graphviz.
%    \begin{macrocode}
\newcommand\pgph@checkrender{%
  \IfFileExists{\pgph@file-proofgraph.\pgph@format}{}%
    {\PackageWarningNoLine{proofgraph}{The autorun option produced no
       \pgph@file-proofgraph.\pgph@format.\MessageBreak
       \ifcase\ShellEscapeStatus
         Shell escape is disabled: compile with -shell-escape%
       \or
         Check that \pgph@engine\space is installed and on the PATH%
       \else
         Shell escape is restricted, and does not permit
         \pgph@engine;\MessageBreak compile with -shell-escape%
       \fi}}%
}
\newcommand\pgph@maybrender{%
  \edef\pgph@sig{\ifdefined\pdf@filemdfivesum
    \pdf@filemdfivesum{\pgph@file.dot}\fi}%
  \ifx\pgph@sig\@empty
    \pgph@render
  \else
    \edef\pgph@sig{\pgph@sig-\pgph@engine-\pgph@format}%
    \edef\pgph@sig{\expandafter\detokenize\expandafter{\pgph@sig}}%
    \IfFileExists{\pgph@file-proofgraph.\pgph@format}%
      {\IfFileExists{\pgph@file-proofgraph.plain}%
        {\ifx\pgph@sig\pgph@prevsig\else\pgph@render\fi}%
        {\pgph@render}}%
      {\pgph@render}%
%    \end{macrocode}
% The signature is recorded whether or not we have just rendered: either way the
% output files on disk are the ones this |.dot| produces. The write must be
% |\immediate|, as |\pgph@emit| runs at end of document, where a deferred write
% would have no shipout left to carry it.
%    \begin{macrocode}
    \if@filesw
      \immediate\write\@auxout{\string\pgph@lastrender{\pgph@sig}}%
    \fi
  \fi
}
%    \end{macrocode}
% \end{macro}
% \end{macro}
% \end{macro}
% \end{macro}
%
% \begin{macro}{\pgph@emit}
% Open the file, resolve editing rules, write nodes then edges, close,
% and optionally run \Graphviz.
%    \begin{macrocode}
\newwrite\pgph@out
\newwrite\pgph@mapout
\newcommand\pgph@emit{%
  \immediate\openout\pgph@out=\pgph@file.dot\relax
  \immediate\openout\pgph@mapout=\pgph@file-proofgraph.map\relax
  \immediate\write\pgph@out{digraph proofgraph \pgph@ob}%
  \immediate\write\pgph@out{\space\space rankdir="\pgph@rankdir";}%
  \begingroup
    \pgph@ignores
    \pgph@excludes
    \pgph@nodes
    \pgph@edges
    \pgph@checkrules
  \endgroup
  \immediate\write\pgph@out{\pgph@cb}%
  \immediate\closeout\pgph@out
  \immediate\closeout\pgph@mapout
  \ifpgph@autorun\pgph@maybrender\fi
}
%    \end{macrocode}
% The emission is registered at |\begin{document}| rather than at load time
% so that it is added to the end-document hook \emph{after} packages loaded
% later (notably \apxproof, whose appendix, with the deferred proofs and
% their references, is built at end of document). This way the graph is
% written only once that material has been typeset, and the emission does
% not disturb the verbatim replay of deferred proofs.
%    \begin{macrocode}
\AtBeginDocument{\AtEndDocument{\pgph@emit}}
%    \end{macrocode}
% If \textsf{TikZ} is available, load its |calc| library now (at
% |\begin{document}|, where category codes are normal) so that the hyperlink
% overlay in |\proofgraph| can use it without reading a file mid-document.
%    \begin{macrocode}
\AtBeginDocument{\@ifpackageloaded{tikz}{\usetikzlibrary{calc}}{}}
%    \end{macrocode}
% \end{macro}
%
% \begin{macro}{\pgphnodemap}
% \begin{macro}{\pgphnodemapcite}
% One entry of the node-to-target map written by |\pgph@emit|: it records, for a
% \Graphviz\ node identifier, what the node should link to. |\pgphnodemap| records
% the user \emph{label} of the result a node represents; |\pgphnodemapcite| records
% the \emph{anchor} of an external citation node, so that (when the cited
% work is in the same document) the node can link to its bibliography entry,
% \textsf{hyperref}'s |cite.|\meta{anchor} destination, the anchor being the
% \BibTeX\ key except under \textsf{biblatex} (see |\pgph@citeanchor|). The map
% is read back when the graph is embedded, to turn each node into a hyperlink.
%    \begin{macrocode}
\newcommand\pgphnodemap[2]{\global\@namedef{pgph@lbl@#1}{#2}}
\newcommand\pgphnodemapcite[2]{\global\@namedef{pgph@clbl@#1}{#2}}
%    \end{macrocode}
% \end{macro}
% \end{macro}
%
% \begin{macro}{\proofgraph}
% Include the rendered graph (from a previous run) in |autorun| mode. When
% \textsf{hyperref} and \textsf{TikZ} are both available (and |hyperlinks| is not
% switched off), each node is made into an internal hyperlink to what it
% represents: a result node links to the result (its |\label|), and an external
% citation node (from the |cite| option) links to the cited work's bibliography
% entry, the |cite.|\meta{anchor} destination \textsf{hyperref} gave it, when
% that work is cited in the same document.
% \Graphviz\ produces a |.pdf| whose link annotations
% would be lost by |\includegraphics| (they sit inside a form \textsc{xobject}),
% so instead we overlay transparent link areas at the node positions, which we
% read from the |-Tplain| output. Without those packages we fall back to a plain
% inclusion.
%
% The image not being there is only to be expected under |autorun| on a first
% run, where it is produced at the end of the run and embedded by the next one,
% so the warning says so rather than pointing at shell escape and the option,
% which the reader has then already seen to. What is worth reporting when
% \Graphviz\ has not run is reported where that is known, by |\pgph@checkrender|
% at the end of the run.
%    \begin{macrocode}
\newif\ifpgph@canlink
\newif\ifpgph@link
\newsavebox\pgph@imgbox
\newsavebox\pgph@natbox
\newread\pgph@plainin
\def\pgph@stop{}
\newcommand\proofgraph[1][]{%
  \@ifundefined{includegraphics}%
    {\PackageError{proofgraph}{\string\proofgraph\space requires the
       graphicx package}{Add \string\usepackage{graphicx} to your preamble
       to embed the rendered graph.}}%
    {\IfFileExists{\pgph@file-proofgraph.\pgph@format}%
      {\pgph@includegraph{#1}}%
      {\PackageWarning{proofgraph}{Rendered graph
         \pgph@file-proofgraph.\pgph@format\space not found;
         \ifpgph@autorun it is rendered at the end of this run, so
         re-run to embed it\else compile with shell-escape and the
         autorun option, then re-run\fi}}}%
}
\newcommand\pgph@includegraph[1]{%
  \pgph@canlinkfalse
  \ifpgph@hyperlinks
    \@ifpackageloaded{hyperref}%
      {\@ifpackageloaded{tikz}%
         {\IfFileExists{\pgph@file-proofgraph.plain}%
            {\pgph@canlinktrue}{}}{}}{}%
  \fi
  \ifpgph@canlink
    \pgph@graphlinked{#1}%
  \else
    \includegraphics[#1]{\pgph@file-proofgraph.\pgph@format}%
  \fi
}
%    \end{macrocode}
% The linked version. We read the node-to-target map under safe category codes (a
% label may contain |_| or, under \textsf{babel}, an active |:|), measure the
% rendered image, place it in a \textsf{TikZ} node, then read the |-Tplain| file
% and drop one transparent link box on each mapped node. \Graphviz\ draws the
% graph inset by a margin, so the image is larger than the layout bounding box
% the |-Tplain| coordinates live in; we recover that margin by also measuring the
% image at its \emph{natural} size and comparing it with the |-Tplain| graph size
% (the |graph| line gives it in inches, hence the factor~72). Node positions are
% then expressed as fractions of the natural image, so they line up with the
% drawn nodes whatever size the image is finally scaled to. The |node| lines
% echo the style attributes of each node, so a colour such as |"#ffd700"| set
% with |\proofgraphstyle| comes back at us; the file is therefore read with
% |#| as an ordinary character, which the \textsf{tikzpicture} group keeps
% local.
%    \begin{macrocode}
\newcommand\pgph@graphlinked[1]{%
  \begingroup
    \catcode`\:=12 \catcode`\_=12 \catcode`\;=12 \catcode`\!=12
    \catcode`\?=12 \catcode`\&=12 \catcode`\^=12 \catcode`\~=12
    \InputIfFileExists{\pgph@file-proofgraph.map}{}{}%
  \endgroup
  \sbox\pgph@imgbox{%
    \includegraphics[#1]{\pgph@file-proofgraph.\pgph@format}}%
  \sbox\pgph@natbox{%
    \includegraphics{\pgph@file-proofgraph.\pgph@format}}%
  \edef\pgph@imgw{\the\wd\pgph@imgbox}%
  \edef\pgph@imgh{\the\dimexpr\ht\pgph@imgbox+\dp\pgph@imgbox\relax}%
  \edef\pgph@natw{\the\wd\pgph@natbox}%
  \edef\pgph@nath{\the\dimexpr\ht\pgph@natbox+\dp\pgph@natbox\relax}%
  \begin{tikzpicture}
    \node[inner sep=\z@,outer sep=\z@](pgph@P){\usebox\pgph@imgbox};
    \catcode`\#=12
    \openin\pgph@plainin=\pgph@file-proofgraph.plain\relax
    \pgph@loopplain
    \closein\pgph@plainin
  \end{tikzpicture}%
}
\newcommand\pgph@loopplain{%
  \unless\ifeof\pgph@plainin
    \read\pgph@plainin to \pgph@line
    \pgph@procline
    \expandafter\pgph@loopplain
  \fi
}
\newcommand\pgph@procline{%
  \IfBeginWith{\pgph@line}{graph }%
    {\expandafter\pgph@dograph\pgph@line\pgph@stop}%
    {\IfBeginWith{\pgph@line}{node }%
       {\expandafter\pgph@donode\pgph@line\pgph@stop}{}}%
}
\def\pgph@dograph graph #1 #2 #3 #4\pgph@stop{%
  \def\pgph@gw{#2}\def\pgph@gh{#3}}
\def\pgph@donode node #1 #2 #3 #4 #5 #6\pgph@stop{%
  \StrDel{#1}{"}[\pgph@nid]%
  \pgph@linkfalse
  \ifcsname pgph@lbl@\pgph@nid\endcsname
    \edef\pgph@tgt{\csname pgph@lbl@\pgph@nid\endcsname}%
    \edef\pgph@thislink{\noexpand\hyperref[\pgph@tgt]}%
    \pgph@linktrue
  \else
    \ifcsname pgph@clbl@\pgph@nid\endcsname
      \edef\pgph@tgt{\csname pgph@clbl@\pgph@nid\endcsname}%
      \edef\pgph@thislink{\noexpand\hyperlink{cite.\pgph@tgt}}%
      \pgph@linktrue
    \fi
  \fi
  \ifpgph@link
    % Per-side Graphviz margin: half the natural-minus-layout size.
    \pgfmathsetmacro\pgph@mx{(\pgph@natw-\pgph@gw*72)/2}%
    \pgfmathsetmacro\pgph@my{(\pgph@nath-\pgph@gh*72)/2}%
    % Node centre as a fraction of the natural image.
    \pgfmathsetmacro\pgph@fx{(\pgph@mx+#2*72)/\pgph@natw}%
    \pgfmathsetmacro\pgph@fy{(\pgph@my+#3*72)/\pgph@nath}%
    \pgfmathsetlengthmacro\pgph@bw{#4*72/\pgph@natw*\pgph@imgw*0.92}%
    \pgfmathsetlengthmacro\pgph@bh{#5*72/\pgph@nath*\pgph@imgh*0.9}%
    \node[anchor=center]
      at ($($(pgph@P.south west)!\pgph@fx!(pgph@P.south east)$)%
      !\pgph@fy!($(pgph@P.north west)!\pgph@fx!(pgph@P.north east)$)$)
      {\pgph@thislink{\phantom{\rule{\pgph@bw}{\pgph@bh}}}};%
  \fi
}
%    \end{macrocode}
% \end{macro}
%
%    \begin{macrocode}
%</package>
%    \end{macrocode}
%
% \Finale
\endinput
